\documentclass[11pt,a4paper]{article}
\usepackage[UTF8]{ctex}
\usepackage{amsmath,amsthm,amssymb}
\usepackage{mathtools}
\usepackage{enumitem}
\usepackage{geometry}
\usepackage{tikz}
\usepackage{pgfplots}
\usepackage{tcolorbox}
\usepackage{algorithm}
\usepackage{algpseudocode}
\usepackage{booktabs}
\usepackage{hyperref}
\usepackage{xcolor}
\usepackage{framed}
\usepackage{pifont}

\geometry{margin=2.5cm}
\pgfplotsset{compat=1.17}

% 定理环境
\theoremstyle{definition}
\newtheorem{definition}{定义}[section]
\newtheorem{example}{例}[section]
\newtheorem{theorem}{定理}[section]
\newtheorem{lemma}{引理}[section]
\newtheorem{corollary}{推论}[section]
\newtheorem{remark}{注记}[section]
\newtheorem{proposition}{命题}[section]

% 自定义盒子
\tcbuselibrary{skins,breakable}
\newtcolorbox{intuition}[1][直觉]{
  colback=blue!5!white,
  colframe=blue!75!black,
  fonttitle=\bfseries,
  title={#1}
}

\newtcolorbox{keypoint}[1][关键点]{
  colback=red!5!white,
  colframe=red!75!black,
  fonttitle=\bfseries,
  title={#1}
}

\newtcolorbox{plainwords}[1][大白话]{
  colback=yellow!8!white,
  colframe=yellow!60!black,
  fonttitle=\bfseries,
  title={#1}
}

\newtcolorbox{proofidea}[1][证明思路]{
  colback=green!5!white,
  colframe=green!50!black,
  fonttitle=\bfseries,
  title={#1}
}

\newtcolorbox{magicbox}[1][神奇结果]{
  colback=purple!5!white,
  colframe=purple!75!black,
  fonttitle=\bfseries,
  title={#1}
}

\newtcolorbox{calculation}[1][动手算一算]{
  colback=orange!5!white,
  colframe=orange!75!black,
  fonttitle=\bfseries,
  title={#1}
}

\title{\textbf{第12周：MDP与Tabular RL理论} \\[0.5em]
\large 从Bandit到序贯决策，收敛性理论的完整图景}
\author{优化算法课程讲义}
\date{}

\begin{document}
\maketitle

\begin{abstract}
上周我们学习了Bandit问题：每轮选一个arm，拿到奖励，选择之间\textbf{相互独立}。本周我们加入\textbf{状态}：你的选择不仅得到奖励，还会改变世界状态，影响未来的选择。这就是马尔可夫决策过程(MDP)。我们将证明Bellman方程为何成立、值迭代为何收敛，以及各种RL算法的收敛性理论。
\end{abstract}

\tableofcontents
\newpage

%=============================================================================
\section{从Bandit到MDP：一步之遥}
%=============================================================================

\begin{plainwords}[回顾上周]
上周的Bandit问题：
\begin{itemize}
    \item $K$台老虎机，每台有未知的期望奖励$\mu_a$
    \item 每轮选一台拉，观察奖励
    \item 选择之间\textbf{完全独立}——你选什么不影响下一轮的世界
\end{itemize}

我们学了UCB（乐观原则）和Thompson Sampling（概率匹配），Regret都是$O(\sqrt{T})$。
\end{plainwords}

\begin{plainwords}[本周的扩展]
现实中，决策往往\textbf{改变世界状态}：
\begin{itemize}
    \item 下棋：每步棋改变棋盘局面
    \item 开车：每个动作改变位置和速度  
    \item 治疗：每次用药改变病人状态
\end{itemize}

MDP就是在Bandit基础上加一个东西：\textbf{状态转移}。

\begin{center}
\begin{tabular}{l|c|c}
\toprule
& Bandit & MDP \\
\midrule
状态 & 只有1个（或者说没有） & 多个状态$s \in S$ \\
动作效果 & 只影响当前奖励 & 影响奖励 + 转移到新状态 \\
决策 & 独立的 & 序贯的，相互影响 \\
\bottomrule
\end{tabular}
\end{center}
\end{plainwords}

\begin{example}[从老虎机到GridWorld]
\textbf{Bandit}：3台老虎机，选一台拉，得奖励，结束。下一轮还是面对同样的3台机器。

\textbf{MDP (GridWorld)}：你在一个$4\times4$网格里，每个格子是一个\textbf{状态}。

\begin{center}
\begin{tikzpicture}[scale=0.7]
\draw[step=1cm,gray,very thin] (0,0) grid (4,4);
\node at (0.5,0.5) {\small 起点};
\node at (3.5,3.5) {\small +10};
\node at (1.5,2.5) {\small -5};
\draw[thick,->] (0.5,0.7) -- (0.5,1.3);
\draw[thick,->] (0.7,1.5) -- (1.3,1.5);
\end{tikzpicture}
\end{center}

\begin{itemize}
    \item 你在某个格子（状态$s$）
    \item 选择一个方向走（动作$a$）
    \item 可能滑到旁边（随机转移$P(s'|s,a)$）
    \item 到达新格子，得到那个格子的奖励（$R(s,a)$）
    \item 现在你在\textbf{新状态}$s'$，面对的选择变了！
\end{itemize}

这就是Bandit和MDP的核心区别：\textbf{你的选择改变了你面对的世界}。
\end{example}

%=============================================================================
\section{MDP形式化定义}
%=============================================================================

\begin{definition}[马尔可夫决策过程 (MDP)]
MDP是一个五元组$(S, A, P, R, \gamma)$：
\begin{itemize}
    \item $S$：状态空间（$|S|$个状态）
    \item $A$：动作空间（$|A|$个动作）
    \item $P: S \times A \times S \to [0,1]$：转移概率，$P(s'|s,a)$
    \item $R: S \times A \to \mathbb{R}$：奖励函数
    \item $\gamma \in [0,1)$：折扣因子
\end{itemize}
\end{definition}

\begin{example}[咖啡店选址——最简单的MDP]
你要决定每天去哪家咖啡店。

\textbf{状态}$S$：你现在的心情 $\{$开心, 一般, 烦躁$\}$

\textbf{动作}$A$：去哪家店 $\{$星巴克, 瑞幸$\}$

\textbf{转移}$P$：
\begin{itemize}
    \item 心情一般 + 去星巴克（贵但好喝）→ 70\%变开心，30\%还是一般
    \item 心情一般 + 去瑞幸（便宜但一般）→ 40\%变开心，60\%还是一般
    \item 心情烦躁 + 去星巴克 → 50\%变一般，50\%还是烦躁
\end{itemize}

\textbf{奖励}$R$：开心时+2，一般时+1，烦躁时-1

\textbf{问题}：长期来看，应该怎么选咖啡店？
\end{example}

\begin{plainwords}[马尔可夫性质]
"马尔可夫"的意思是：\textbf{未来只依赖于现在，不依赖于过去}。

\textbf{生活例子}：天气预报

"明天是否下雨"只取决于"今天的天气状况"，不取决于"上周下了几天雨"。

如果今天是晴天，明天30\%下雨。不管之前连续晴了3天还是30天，这个概率都是30\%。

这大大简化了问题：我们只需要记住"现在在哪"，不需要记住"怎么来的"。
\end{plainwords}

\begin{example}[为什么需要折扣因子$\gamma$]
假设你可以选择：
\begin{itemize}
    \item 选项A：今天拿100块
    \item 选项B：一年后拿100块
\end{itemize}

大多数人选A！为什么？

\begin{enumerate}
    \item \textbf{不确定性}：一年后的事谁知道呢
    \item \textbf{机会成本}：今天的100块可以投资
    \item \textbf{即时满足}：人类天生更看重眼前
\end{enumerate}

$\gamma$就是"未来的钱打几折"。$\gamma = 0.9$意味着：
\begin{itemize}
    \item 明天的1块 = 今天的0.9块
    \item 后天的1块 = 今天的$0.9^2 = 0.81$块
    \item 一年后的1块 = 今天的$0.9^{365} \approx 0$块
\end{itemize}
\end{example}

\begin{definition}[策略与值函数]
\textbf{策略}$\pi: S \to A$（或$\pi: S \to \Delta(A)$）告诉我们在每个状态选什么动作。

\textbf{值函数}$V^\pi(s)$：从状态$s$出发，按策略$\pi$行动，获得的期望累积折扣奖励：
\[
V^\pi(s) = \mathbb{E}_\pi\left[\sum_{t=0}^\infty \gamma^t R(s_t, a_t) \mid s_0 = s\right]
\]

\textbf{动作-值函数}$Q^\pi(s,a)$：从状态$s$采取动作$a$，然后按$\pi$行动：
\[
Q^\pi(s,a) = R(s,a) + \gamma \sum_{s'} P(s'|s,a) V^\pi(s')
\]
\end{definition}

\begin{example}[值函数的直觉——房产估值]
把$V(s)$想象成\textbf{房子的价值}：

\begin{itemize}
    \item 状态$s$ = 房子的位置
    \item 动作$a$ = 装修方案
    \item 奖励$R$ = 每月租金收入
    \item 转移$P$ = 装修后房子状态变化（可能升值也可能出问题）
\end{itemize}

$V(s)$就是"这个位置的房子，长期能给我带来多少租金收入（折现到今天）"。

好位置（市中心）$V$高，差位置（偏远郊区）$V$低。
\end{example}

\begin{calculation}[$V$和$Q$的关系]
$V^\pi(s) = Q^\pi(s, \pi(s))$（确定性策略）

\textbf{生活比喻}：
\begin{itemize}
    \item $V(s)$ = 在状态$s$，按我的习惯做事，能得多少分
    \item $Q(s,a)$ = 在状态$s$，如果这次选动作$a$，能得多少分
    \item $V(s) = Q(s, \text{我的习惯选择})$
\end{itemize}

\textbf{数值例子}：状态$s$，动作$a$，奖励$R(s,a)=1$，$\gamma=0.9$。

转移：50\%到$s_1$（$V^\pi(s_1)=5$），50\%到$s_2$（$V^\pi(s_2)=3$）。

\[
Q^\pi(s,a) = \underbrace{1}_{\text{立即奖励}} + \underbrace{0.9}_{\text{折扣}} \times \underbrace{(0.5 \times 5 + 0.5 \times 3)}_{\text{未来期望价值}} = 1 + 3.6 = 4.6
\]
\end{calculation}

%=============================================================================
\section{Bellman方程：递归的威力}
%=============================================================================

\begin{plainwords}
Bellman方程是RL的\textbf{牛顿定律}，一切都从它推导出来。

核心思想极其简单：\textbf{今天的价值 = 今天的奖励 + 明天的价值（打折）}

这个递归关系把无限期问题变成了一步一步的计算。
\end{plainwords}

\begin{example}[Bellman方程的生活版——职业规划]
你在考虑职业发展，当前状态是"初级工程师"。

\textbf{问}：这个状态的"价值"是多少？（未来能赚多少钱）

\textbf{Bellman式思考}：
\[
V(\text{初级}) = \underbrace{15\text{万}}_{\text{今年工资}} + 0.9 \times \underbrace{[\text{明年状态的价值}]}_{\text{可能升职/跳槽/原地}}
\]

如果30\%概率升职（高级值50万），70\%概率还是初级：
\[
V(\text{初级}) = 15 + 0.9 \times (0.3 \times 50 + 0.7 \times V(\text{初级}))
\]

解方程：$V = 15 + 13.5 + 0.63V$，得$V(\text{初级}) = 77$万。

\textbf{这就是Bellman方程}：当前价值 = 即时回报 + 折扣后的未来期望！
\end{example}

\subsection{Bellman期望方程}

\begin{theorem}[Bellman期望方程]
对任意策略$\pi$：
\[
V^\pi(s) = \sum_a \pi(a|s) \left[R(s,a) + \gamma \sum_{s'} P(s'|s,a) V^\pi(s')\right]
\]
用算子形式：$V^\pi = T^\pi V^\pi$，即$V^\pi$是Bellman算子$T^\pi$的\textbf{不动点}。
\end{theorem}

\begin{example}[2状态MDP——学生的一天]
\textbf{状态}：\{学习中, 摸鱼中\}

\textbf{转移}（无动作，自动转移）：
\begin{itemize}
    \item 学习中 → 60\%继续学习，40\%开始摸鱼
    \item 摸鱼中 → 30\%开始学习，70\%继续摸鱼
\end{itemize}

\textbf{奖励}：学习+2分，摸鱼+1分。\textbf{$\gamma = 0.9$}

\textbf{Bellman方程}：
\begin{align*}
V(\text{学习}) &= 2 + 0.9 \times (0.6 \times V(\text{学习}) + 0.4 \times V(\text{摸鱼})) \\
V(\text{摸鱼}) &= 1 + 0.9 \times (0.3 \times V(\text{学习}) + 0.7 \times V(\text{摸鱼}))
\end{align*}

\textbf{求解}：从第二个方程整理：
$V(\text{摸鱼}) = 1 + 0.27 V(\text{学习}) + 0.63 V(\text{摸鱼})$

$0.37 V(\text{摸鱼}) = 1 + 0.27 V(\text{学习})$ → $V(\text{摸鱼}) = 2.7 + 0.73 V(\text{学习})$

代入第一个：$V(\text{学习}) = 2 + 0.54 V(\text{学习}) + 0.36(2.7 + 0.73 V(\text{学习}))$

$V(\text{学习}) = 2.97 + 0.8 V(\text{学习})$ → $V(\text{学习}) = 14.85$

$V(\text{摸鱼}) = 2.7 + 0.73 \times 14.85 = 13.54$

\textbf{结论}：学习状态更"值钱"！（14.85 vs 13.54）
\end{example}

\begin{proof}[Bellman期望方程的证明]
从定义展开：
\begin{align*}
V^\pi(s) &= \mathbb{E}_\pi\left[\sum_{t=0}^\infty \gamma^t R(s_t, a_t) \mid s_0 = s\right] \\
&= \mathbb{E}_\pi[R(s_0, a_0) | s_0 = s] + \gamma \mathbb{E}_\pi\left[\sum_{t=1}^\infty \gamma^{t-1} R(s_t, a_t) \mid s_0 = s\right]
\end{align*}

第一项：$\mathbb{E}_\pi[R(s_0, a_0) | s_0 = s] = \sum_a \pi(a|s) R(s,a)$

第二项：利用马尔可夫性，给定$s_1$，未来与$s_0$独立：
\[
\mathbb{E}_\pi\left[\sum_{t=1}^\infty \gamma^{t-1} R(s_t, a_t) \mid s_0 = s\right] = \sum_a \pi(a|s) \sum_{s'} P(s'|s,a) V^\pi(s')
\]

合并得证。
\end{proof}

\subsection{Bellman最优方程}

\begin{definition}[最优值函数]
$V^*(s) = \max_\pi V^\pi(s)$，$Q^*(s,a) = \max_\pi Q^\pi(s,a)$
\end{definition}

\begin{theorem}[Bellman最优方程]
\[
V^*(s) = \max_a \left[R(s,a) + \gamma \sum_{s'} P(s'|s,a) V^*(s')\right]
\]
\[
Q^*(s,a) = R(s,a) + \gamma \sum_{s'} P(s'|s,a) \max_{a'} Q^*(s', a')
\]
\end{theorem}

\begin{example}[期望 vs 最优——午餐选择]
你饿了，两个选择：
\begin{itemize}
    \item 动作A：食堂（确定7分满意度）
    \item 动作B：新餐厅（50\%概率9分超好吃，50\%概率5分踩雷）
\end{itemize}

\textbf{如果你的策略是"随机选"}（各50\%）：
\[
V^{\pi} = 0.5 \times 7 + 0.5 \times (0.5 \times 9 + 0.5 \times 5) = 0.5 \times 7 + 0.5 \times 7 = 7
\]

\textbf{最优策略}：比较两个动作的期望
\begin{itemize}
    \item $Q^*(s, \text{食堂}) = 7$
    \item $Q^*(s, \text{新餐厅}) = 0.5 \times 9 + 0.5 \times 5 = 7$
\end{itemize}

两个一样好！$V^* = 7$。

\textbf{但如果}新餐厅是60\%概率9分：$Q^*(s, \text{新餐厅}) = 7.4 > 7$

那最优策略就是\textbf{总是选新餐厅}，$V^* = 7.4$。
\end{example}

\begin{plainwords}[期望 vs 最优]
\textbf{Bellman期望方程}：按固定策略$\pi$行动的价值，用$\sum_a \pi(a|s)$加权

\textbf{Bellman最优方程}：每一步都选最好的，用$\max_a$

最优策略就是贪婪策略：$\pi^*(s) = \arg\max_a Q^*(s,a)$
\end{plainwords}

%=============================================================================
\section{值迭代：收缩映射的魔力}
%=============================================================================

\begin{plainwords}
知道Bellman方程是不动点后，怎么\textbf{找到}它？

最简单的方法：\textbf{迭代}！从任意初始值开始，反复应用Bellman更新。

为什么能行？因为Bellman算子是\textbf{收缩映射}——每次迭代都让你更接近真解！
\end{plainwords}

\begin{example}[收缩映射的生活比喻——调节水温]
你想把浴缸水温调到40°C，但你不知道现在多少度。

\textbf{策略}：每次加一点热水或冷水，让温度往40°C靠近。

如果你的调节方式是"每次让温度往40°C移动一半"：
\begin{itemize}
    \item 现在60°C → 加冷水 → 变成50°C（离目标近了一半）
    \item 现在50°C → 继续调 → 变成45°C
    \item 现在45°C → 继续调 → 变成42.5°C
    \item ...
\end{itemize}

\textbf{关键性质}：每次操作后，误差\textbf{至少减半}！

这就是"收缩映射"：每次操作让你离目标更近，而且有\textbf{保证的进步}。

值迭代就是这样：每次Bellman更新，误差至少乘以$\gamma$（比如0.9）。
\end{example}

\subsection{收缩映射与Banach不动点定理}

\begin{definition}[收缩映射]
算子$T$是$\gamma$-收缩映射，如果对所有$V, V'$：
\[
\|TV - TV'\|_\infty \leq \gamma \|V - V'\|_\infty
\]
\end{definition}

\begin{example}[理解收缩映射——橡皮筋]
想象一根橡皮筋上有两个点$V$和$V'$，距离是$d$。

"收缩映射"就像\textbf{同时拉两个点}，但拉完之后距离变成$\gamma \cdot d$。

如果$\gamma = 0.9$，每次操作距离缩短10\%。
如果$\gamma = 0.5$，每次操作距离缩短50\%。

不管一开始两点离多远，反复操作后它们会\textbf{越来越近}，最终重合！

重合的那个点就是\textbf{不动点}——所有起点最终都会到达那里。
\end{example}

\begin{theorem}[Bellman算子是收缩映射]
Bellman最优算子$(T^* V)(s) = \max_a [R(s,a) + \gamma \sum_{s'} P(s'|s,a) V(s')]$是$\gamma$-收缩。
\end{theorem}

\begin{proof}
对任意$V, V'$和状态$s$，设$a^* = \arg\max_a [R(s,a) + \gamma \sum_{s'} P(s'|s,a) V(s')]$。

\begin{align*}
(T^* V)(s) - (T^* V')(s) &\leq [R(s,a^*) + \gamma \sum_{s'} P(s'|s,a^*) V(s')] - [R(s,a^*) + \gamma \sum_{s'} P(s'|s,a^*) V'(s')] \\
&= \gamma \sum_{s'} P(s'|s,a^*) [V(s') - V'(s')] \\
&\leq \gamma \|V - V'\|_\infty
\end{align*}

类似地可证$(T^* V')(s) - (T^* V)(s) \leq \gamma \|V - V'\|_\infty$。

取max得$\|T^* V - T^* V'\|_\infty \leq \gamma \|V - V'\|_\infty$。
\end{proof}

\begin{plainwords}[证明的直觉]
为什么Bellman算子是收缩的？

因为Bellman更新里有个$\gamma$在前面！

$V' = R + \gamma \times (\text{未来价值})$

未来价值的误差被$\gamma$打了折扣，所以误差在传播过程中\textbf{自然衰减}。
\end{plainwords}

\begin{theorem}[Banach不动点定理]
若$T$是完备度量空间上的$\gamma$-收缩映射（$\gamma < 1$），则：
\begin{enumerate}
    \item $T$有\textbf{唯一}不动点$V^*$
    \item 从任意$V_0$出发，$V_k = T^k V_0 \to V^*$
    \item 收敛速度：$\|V_k - V^*\|_\infty \leq \gamma^k \|V_0 - V^*\|_\infty$
\end{enumerate}
\end{theorem}

\begin{example}[Banach定理的直觉——GPS导航]
你要去一个目的地（不动点$V^*$），但你不知道在哪。

\textbf{GPS的承诺}：每走一步，你离目的地的距离\textbf{至少缩短10\%}。

那么：
\begin{itemize}
    \item 不管你从哪出发，最终一定能到
    \item 只有一个目的地（唯一性）
    \item 走$k$步后，距离最多是初始的$0.9^k$
\end{itemize}

值迭代就是这样的GPS：每次Bellman更新保证让你更接近$V^*$。
\end{example}

\begin{proof}[Banach不动点定理的证明]
\textbf{存在性}：取任意$V_0$，令$V_k = T^k V_0$。对$m > n$：
\begin{align*}
\|V_m - V_n\| &\leq \sum_{i=n}^{m-1} \|V_{i+1} - V_i\| \leq \sum_{i=n}^{m-1} \gamma^i \|V_1 - V_0\| \leq \frac{\gamma^n}{1-\gamma} \|V_1 - V_0\|
\end{align*}

当$n \to \infty$，右边趋于0，所以$\{V_k\}$是Cauchy序列，收敛到某个$V^*$。

由$T$的连续性，$TV^* = T(\lim_k V_k) = \lim_k TV_k = \lim_k V_{k+1} = V^*$。

\textbf{唯一性}：若$V^*, \tilde{V}^*$都是不动点，则：
\[
\|V^* - \tilde{V}^*\| = \|TV^* - T\tilde{V}^*\| \leq \gamma \|V^* - \tilde{V}^*\|
\]
只能$V^* = \tilde{V}^*$。
\end{proof}

\begin{algorithm}[H]
\caption{值迭代 (Value Iteration)}
\begin{algorithmic}[1]
\State 初始化$V_0(s) = 0$对所有$s$
\For{$k = 0, 1, 2, \ldots$}
    \For{每个状态$s$}
        \State $V_{k+1}(s) = \max_a \left[R(s,a) + \gamma \sum_{s'} P(s'|s,a) V_k(s')\right]$
    \EndFor
    \If{$\|V_{k+1} - V_k\|_\infty < \epsilon(1-\gamma)/(2\gamma)$}
        \State 停止
    \EndIf
\EndFor
\State 输出$\pi(s) = \arg\max_a [R(s,a) + \gamma \sum_{s'} P(s'|s,a) V_k(s')]$
\end{algorithmic}
\end{algorithm}

\begin{calculation}[收敛速度]
$\gamma = 0.9$，初始误差$\|V_0 - V^*\| = 100$。

\begin{center}
\begin{tabular}{c|c|c}
\toprule
迭代次数$k$ & $\gamma^k$ & 误差上界 \\
\midrule
10 & 0.35 & 35 \\
50 & 0.005 & 0.5 \\
100 & $3\times10^{-5}$ & 0.003 \\
\bottomrule
\end{tabular}
\end{center}

误差每$\log_{1/\gamma} 10 = \frac{\ln 10}{\ln(1/0.9)} \approx 22$次迭代减少10倍。

\textbf{$\gamma$越大，收敛越慢}：$\gamma=0.99$需要230次迭代误差减10倍。
\end{calculation}

%=============================================================================
\section{策略迭代：更快的收敛}
%=============================================================================

\begin{plainwords}
值迭代每次只更新一点点。策略迭代更激进：

\begin{enumerate}
    \item \textbf{策略评估}：完全解出当前策略$\pi$的值函数$V^\pi$
    \item \textbf{策略改进}：根据$V^\pi$贪婪地更新策略
\end{enumerate}

每次跳到"当前策略下的最优"，而不是一点点挪。
\end{plainwords}

\begin{theorem}[策略改进定理]
设$\pi'(s) = \arg\max_a Q^\pi(s,a)$。则$V^{\pi'}(s) \geq V^\pi(s)$对所有$s$成立。

等号成立当且仅当$\pi$已经是最优策略。
\end{theorem}

\begin{proof}
对任意$s$：
\begin{align*}
V^\pi(s) &\leq \max_a Q^\pi(s,a) = Q^\pi(s, \pi'(s)) \\
&= R(s, \pi'(s)) + \gamma \sum_{s'} P(s'|s, \pi'(s)) V^\pi(s') \\
&\leq R(s, \pi'(s)) + \gamma \sum_{s'} P(s'|s, \pi'(s)) \max_{a'} Q^\pi(s', a')
\end{align*}

继续展开，得到$V^\pi(s) \leq V^{\pi'}(s)$。
\end{proof}

\begin{theorem}[策略迭代有限步收敛]
策略迭代最多$|A|^{|S|}$次迭代收敛到最优策略。

实际中通常只需5-20次迭代。
\end{theorem}

%=============================================================================
\section{Monte Carlo方法：最朴素的model-free方法}
%=============================================================================

\begin{plainwords}
值迭代、策略迭代都假设我们\textbf{完全知道}$P, R$。

现实里更常见的场景：你有一个能跑的环境（机器人、游戏、模拟器），可以执行动作、观察反馈，但$P(s'|s,a)$和$R(s,a)$的解析形式不知道。

\textbf{最朴素的对策}：跑很多次，对每个状态出现时的累积回报取样本均值——这就是Monte Carlo方法。

它是bandit最直接的MDP版本：把"拉一次得到$r$"换成"跑完一个episode得到$G_t$"，剩下全靠大数定律。
\end{plainwords}

\subsection{Bandit的直接推广}

\begin{intuition}[把状态当成"轨迹老虎机"]
上周bandit的核心套路：
\begin{itemize}
\item 不知道$\mu_a$
\item 拉臂$a$多次，记录奖励$r$
\item 用样本均值$\hat\mu_a = \frac{1}{N_a}\sum_{i=1}^{N_a} r_i$估计$\mu_a$
\item 大数定律 $\Rightarrow$ $\hat\mu_a \to \mu_a$
\end{itemize}

MDP里把"臂"换成"状态"，"一次奖励"换成"从该状态出发的整段回报$G_t$"：
\begin{itemize}
\item 不知道$V^\pi(s)$
\item 跑$\pi$多次，每次记录从$s$出发观察到的$G_t$
\item 用样本均值$\hat V(s) = \frac{1}{N_s}\sum_{i} G^{(i)}$估计$V^\pi(s)$
\end{itemize}

\textbf{大数定律仍然保证收敛}——因为定义上$V^\pi(s) = \mathbb{E}[G_t \mid s_t = s]$，样本均值估期望嘛。

\textbf{这就是Monte Carlo方法。}它和bandit是同一个统计原理，只不过把"瞬时奖励"换成了"未来累积折扣奖励"。
\end{intuition}

\subsection{MC评估算法（prediction）}

\begin{definition}[首次访问MC评估]
给定策略$\pi$。运行$K$条完整轨迹：
\begin{algorithm}[H]
\caption{First-Visit Monte Carlo评估}
\begin{algorithmic}[1]
\State 初始化$V(s) = 0$，访问计数$N(s) = 0$
\For{episode $k = 1, \ldots, K$}
    \State 用$\pi$生成完整轨迹$(s_0, a_0, r_0, s_1, a_1, r_1, \ldots, s_T)$
    \State 计算$G_t = \sum_{j=t}^{T-1} \gamma^{j-t} r_j$对每个$t$
    \For{每个在该轨迹中\textbf{首次出现}的状态$s$，对应时刻$t$}
        \State $N(s) \mathrel{+}= 1$
        \State $V(s) \leftarrow V(s) + \frac{1}{N(s)}\big[G_t - V(s)\big]$
    \EndFor
\EndFor
\end{algorithmic}
\end{algorithm}

更新规则展开等价于增量样本均值：
\[
V_n(s) = \frac{1}{n}\sum_{i=1}^n G^{(i)}
\]
\end{definition}

\begin{remark}[首次访问 vs 每次访问]
\textbf{First-visit MC}：一条轨迹里$s$出现多次只用第一次的$G$。样本独立，方差分析干净。

\textbf{Every-visit MC}：$s$每次出现都更新。样本量更大但相关。

两者都w.p.1收敛到$V^\pi$。第一种严格无偏，第二种渐近无偏。实践用every-visit的多（数据利用率高）。
\end{remark}

\subsection{MC的统计性质}

\begin{theorem}[MC评估的收敛]
若每个状态$s$被无限次访问，且步长满足$\sum_n \alpha_n = \infty, \sum_n \alpha_n^2 < \infty$（或用$\alpha_n = 1/n$），则
\[
V_n(s) \xrightarrow{w.p.1} V^\pi(s)
\]
\end{theorem}

\begin{proofidea}
这\textbf{不是新结果}。$V^\pi(s) = \mathbb{E}[G_t \mid s_t = s]$是定义，$\{G^{(i)}\}$（首次访问版本）是i.i.d.\ 样本，强大数定律即得。

\textbf{MC评估在统计学上就是bandit的样本均值定理在MDP上的复读机版本。}没有新的概率工具，只是把"奖励"重新定义成"未来累积折扣奖励"。
\end{proofidea}

\begin{calculation}[MC的方差量级]
单次$G_t$的方差，假设每步奖励方差$\sigma^2$、各步独立（最乐观情形）：
\[
\text{Var}(G_t) = \sum_{k=0}^\infty \gamma^{2k}\sigma^2 = \frac{\sigma^2}{1-\gamma^2}
\]
$\gamma=0.9 \Rightarrow$ 放大$\approx 5.3$倍；$\gamma=0.99 \Rightarrow$ 放大$\approx 50$倍。

样本均值的标准差是$\sqrt{\text{Var}(G_t)/n}$，要达到精度$\epsilon$需要：
\[
n = O\Big(\frac{\sigma^2}{(1-\gamma^2)\epsilon^2}\Big)
\]
轨迹。$\gamma$越接近1，常数爆炸。

实际情形（状态转移随机性叠加奖励随机性）方差更大。\textbf{这是MC的根本痛点。}
\end{calculation}

\subsection{MC control：评估 + 改进的循环}

Prediction解决了，control（找最优$\pi$）怎么办？沿用策略迭代的GPI框架：评估 + greedy改进。

\begin{plainwords}[Model-free control的关键切换：必须学$Q$不学$V$]
策略迭代的改进步骤：
\[
\pi'(s) = \arg\max_a Q^\pi(s,a)
\]
如果只有$V^\pi$，model-based下可以算$Q$：
\[
Q^\pi(s,a) = R(s,a) + \gamma\sum_{s'} P(s'|s,a)V^\pi(s')
\]
——\textbf{这需要$P, R$}！

所以：\textbf{model-free control必须直接估计$Q$，不能只估$V$}。这是从prediction（学$V$就够）跳到control（必须学$Q$）的根本切换点。

记住这个原则——下一章的TD也会撞上同一堵墙，得出同样的结论。
\end{plainwords}

\begin{algorithm}[H]
\caption{MC Control（$\epsilon$-greedy版）}
\begin{algorithmic}[1]
\State 初始化$Q(s,a)$任意，$\pi = \epsilon$-greedy w.r.t.\ $Q$
\For{episode $k = 1, 2, \ldots$}
    \State 用$\pi$生成完整轨迹
    \State 对每个访问到的$(s,a)$：$Q(s,a) \leftarrow Q(s,a) + \alpha[G - Q(s,a)]$
    \State $\pi \leftarrow \epsilon$-greedy w.r.t.\ 更新后的$Q$
\EndFor
\end{algorithmic}
\end{algorithm}

\begin{remark}[为什么需要$\epsilon$-greedy]
纯greedy会漏掉好动作——没访问过的$(s,a)$永远没有$Q$估计、永远不会被选中、永远没机会被试。$\epsilon$-greedy保证每个动作都有$\geq \epsilon/|A|$的概率被尝试，对应bandit的\textbf{forced exploration}。

GLIE条件（$\epsilon \to 0$、每个$(s,a)$被无限次访问）下，MC control w.p.1收敛到$Q^*$。
\end{remark}

\subsection{MC的两个软肋}

MC虽然简单优雅、和bandit理论同源，但有两个严重缺陷：

\begin{keypoint}[MC不适用的场景]
\begin{enumerate}
\item \textbf{必须等到episode结束}才能算$G_t$。
\begin{itemize}
\item 无限horizon任务（连续运行的工厂控制、未通关的游戏）根本没法用
\item 即使有限horizon，episode长意味着延迟更新——在线学习效率低
\end{itemize}

\item \textbf{方差极大}：
\begin{itemize}
\item $G_t = r_t + \gamma r_{t+1} + \gamma^2 r_{t+2} + \cdots$是$O(1/(1-\gamma))$步噪声的叠加
\item Bandit拉一次承担一步随机性，MC拉一次承担一整段
\item 学习率很难调——大了不稳，小了慢
\end{itemize}
\end{enumerate}
\end{keypoint}

\begin{plainwords}[两条改进路线 = 后面两章]
MC的两个软肋分别催生了后面两章的主题。

\textbf{软肋1（样本效率低、方差大）$\Rightarrow$ §UCB-VI}：与其暴力 sample-average 整段回报，不如从样本里\textbf{估出模型}$\hat P, \hat R$，再做optimistic planning。把所有样本的信息塞进模型，对每个$(s,a)$都做了"信息共享"——这是\textbf{model-based RL}，能拿到$\tilde O(\sqrt{H^3 SAK})$的regret。

\textbf{软肋2（必须等episode结束）$\Rightarrow$ §时序差分学习}：用\textbf{bootstrap}——一步样本$+$ 自己当前的估计——绕过等待，每步都能更新。保持model-free和online，导出SARSA和Q-learning。

后续这两章请始终回头比较：
\begin{itemize}
\item UCB-VI：sample → 估模型 → 计划（每步$O(SA \cdot S)$存储 + planning开销）
\item TD：sample → 一步bootstrap → 更新$Q$（每步$O(1)$）
\item MC：sample → 等到底 → 取均值（每步$O(1)$但要等episode）
\end{itemize}
三者各自付出不同代价，换不同的好处。
\end{plainwords}

%=============================================================================
\section{从Bandit Regret到RL Regret}
%=============================================================================

\begin{plainwords}[承接上周]
上周定义的Bandit Regret：$\sum_{t=1}^T (\mu^* - \mu_{a_t})$

MDP的Regret类似，但要考虑\textbf{episode结构}：
\end{plainwords}

\begin{definition}[Episodic MDP的Regret]
运行$K$个episode，每个episode长度$H$：
\[
\text{Regret}(K) = \sum_{k=1}^K \left[V_1^*(s_1^k) - V_1^{\pi_k}(s_1^k)\right]
\]
其中$\pi_k$是第$k$个episode使用的策略。
\end{definition}

\begin{keypoint}[Bandit vs MDP Regret]
\begin{center}
\begin{tabular}{l|c|c}
\toprule
& Bandit & MDP \\
\midrule
单轮损失 & $\mu^* - \mu_a$ & $V^*(s) - V^\pi(s)$ \\
累积 & $T$轮 & $K$个episode \\
最优界 & $O(\sqrt{KT})$ & $O(\sqrt{H^3 SAK})$ \\
\bottomrule
\end{tabular}
\end{center}

MDP多了$\sqrt{H^3 SA}$因子：状态空间、动作空间、时间范围都增加了探索难度。
\end{keypoint}

%=============================================================================
\section{UCB-VI：把UCB推广到MDP}
%=============================================================================

\begin{plainwords}[相对MC的改进点]
MC暴力sample-average整段回报，样本效率低。UCB-VI走另一条路：从样本里\textbf{估出模型}$\hat P, \hat R$，再做optimistic planning。
\begin{itemize}
\item MC：每个$(s,a)$的回报\textbf{独立估计}，信息不共享
\item UCB-VI：所有样本\textbf{都塞进模型}$\hat P, \hat R$里，不同$(s,a)$通过planning共享信息
\end{itemize}
代价：要存$O(S^2 A)$的转移计数表 + 每个episode跑一次值迭代。回报：$\tilde O(\sqrt{H^3 SAK})$ regret，远优于MC的样本复杂度。

下面具体看怎么做。
\end{plainwords}

\begin{plainwords}[核心思想]
上周的UCB：$\hat{\mu}_a + \sqrt{\frac{2\ln t}{n_a}}$（均值 + 不确定性bonus）

UCB-VI的思想完全一样：对不确定的\textbf{转移概率}，假设它们对我们\textbf{有利}。

具体做法：在值迭代中，给Q值加一个bonus，bonus大小与$(s,a)$的访问次数成反比。
\end{plainwords}

\begin{example}[UCB-VI的生活版——探索新城市]
你刚搬到一个新城市，想找最好的生活方式。

\textbf{状态}：你在哪个街区

\textbf{动作}：去哪个餐厅/商店/公园

\textbf{问题}：你不知道每个地方有多好，也不知道去了之后会遇到什么（转移概率）。

\textbf{UCB-VI策略}：
\begin{enumerate}
    \item 对每个没去过的地方，\textbf{乐观假设}它很棒（给高分）
    \item 去过几次的地方，根据经验评分，但加一个"不确定性bonus"
    \item 去过很多次的地方，bonus很小，评分接近真实
\end{enumerate}

\textbf{效果}：
\begin{itemize}
    \item 一开始疯狂探索（所有地方bonus都高）
    \item 慢慢bonus降低，开始exploitation
    \item 最终找到最好的生活方式
\end{itemize}

\textbf{和上周UCB的区别}：
\begin{itemize}
    \item UCB：只考虑"这个餐厅有多好"
    \item UCB-VI：还要考虑"去完这个餐厅，我会到哪个状态，那个状态有多好"
\end{itemize}
\end{example}

\begin{algorithm}[H]
\caption{UCB-VI（完整版）}
\begin{algorithmic}[1]
\State 初始化计数$N(s,a) = 0$，$N(s,a,s') = 0$
\For{episode $k = 1, 2, \ldots, K$}
    \State \textbf{// 乐观值迭代（从后往前）}
    \State $V_{H+1}(s) = 0$ for all $s$
    \For{$h = H, H-1, \ldots, 1$}
        \For{each $(s, a)$}
            \If{$N(s,a) > 0$}
                \State $\hat{P}(s'|s,a) = N(s,a,s') / N(s,a)$
                \State $\text{bonus}_h(s,a) = c\sqrt{\frac{H^2 \iota}{N(s,a)}}$ where $\iota = \ln(SAH K/\delta)$
            \Else
                \State $\text{bonus}_h(s,a) = H$ \Comment{最大可能值}
            \EndIf
            \State $Q_h(s,a) = \min\left\{H, R(s,a) + \sum_{s'} \hat{P}(s'|s,a) V_{h+1}(s') + \text{bonus}_h(s,a)\right\}$
        \EndFor
        \State $V_h(s) = \max_a Q_h(s,a)$, \quad $\pi_h^k(s) = \arg\max_a Q_h(s,a)$
    \EndFor
    \State \textbf{// 执行策略并收集数据}
    \State 从$s_1^k$出发，执行$\pi^k$，观察轨迹$(s_1, a_1, r_1, \ldots, s_H, a_H, r_H)$
    \State 更新计数：$N(s_h, a_h) \mathrel{+}= 1$，$N(s_h, a_h, s_{h+1}) \mathrel{+}= 1$
\EndFor
\end{algorithmic}
\end{algorithm}

\subsection{UCB-VI的Regret分析}

\begin{theorem}[UCB-VI的Regret界]
以概率至少 $1-\delta$，有
\[
\operatorname{Regret}(K)
=
\sum_{k=1}^K
\left[
V_1^*(s_1^k)-V_1^{\pi^k}(s_1^k)
\right]
\leq
O\left(\sqrt{H^3SAK\iota}\right),
\]
其中
\[
\iota=\ln(SAHK/\delta).
\]
\end{theorem}

\begin{proof}
证明分四步。

\textbf{Step 1：建立好事件}

定义好事件 $\mathcal{E}$：对所有 $(s,a,h,k)$，经验转移概率的误差均被 bonus 控制，即
\[
\left|
\sum_{s'}
\left(
\hat P^k(s'|s,a)-P(s'|s,a)
\right)
V_{h+1}^k(s')
\right|
\leq
\operatorname{bonus}_h^k(s,a),
\]
其中
\[
\operatorname{bonus}_h^k(s,a)
=
c\sqrt{
\frac{H^2\iota}{N^k(s,a)\vee 1}
}.
\]
这里 \(N^k(s,a)\) 表示在 episode \(k\) 开始前，状态动作对 \((s,a)\) 被访问的总次数。注意这里的计数不依赖于 \(h\)。

\begin{lemma}[好事件成立]
选择足够大的常数 $c$，有
\[
\mathbb{P}(\mathcal{E})\geq 1-\delta.
\]
\end{lemma}

\begin{proof}[引理证明]
固定 $(s,a,h,k)$。设在 episode \(k\) 开始前，\((s,a)\) 已经被访问了
\[
n=N^k(s,a)
\]
次。由于 \(V_{h+1}^k(s')\in[0,H]\)，由 Hoeffding 不等式，
\[
\mathbb{P}
\left(
\left|
\sum_{s'}
(\hat P^k-P)(s'|s,a)
V_{h+1}^k(s')
\right|
>
\epsilon
\right)
\leq
2\exp\left(
-\frac{2n\epsilon^2}{H^2}
\right).
\]
取
\[
\epsilon
=
c\sqrt{\frac{H^2\iota}{n\vee 1}},
\]
则失败概率至多为
\[
2\exp(-2c^2\iota)
\leq
\frac{\delta}{SAHK},
\]
其中最后一步可通过选择足够大的常数 \(c\) 保证。

对所有 $(s,a,h,k)$ 作 union bound，得到
\[
\mathbb{P}(\mathcal{E})\geq 1-\delta.
\]
\end{proof}

\textbf{Step 2：乐观性}

\begin{lemma}[乐观引理]
在好事件 $\mathcal{E}$ 下，对所有 $(s,h,k)$，有
\[
V_h^k(s)\geq V_h^*(s).
\]
\end{lemma}

\begin{proof}[引理证明]
对 $h$ 从 $H+1$ 到 $1$ 作反向归纳。

当 \(h=H+1\) 时，
\[
V_{H+1}^k(s)=0=V_{H+1}^*(s).
\]

假设对 \(h+1\) 有
\[
V_{h+1}^k(s)\geq V_{h+1}^*(s).
\]
则对任意 $(s,a)$，
\begin{align*}
Q_h^k(s,a)
&=
R(s,a)
+
\sum_{s'}\hat P^k(s'|s,a)V_{h+1}^k(s')
+
\operatorname{bonus}_h^k(s,a)
\\
&\geq
R(s,a)
+
\sum_{s'}P(s'|s,a)V_{h+1}^k(s')
\\
&\geq
R(s,a)
+
\sum_{s'}P(s'|s,a)V_{h+1}^*(s')
\\
&=
Q_h^*(s,a).
\end{align*}
第一处不等式使用了好事件，第二处不等式使用了归纳假设。

因此
\[
V_h^k(s)
=
\max_a Q_h^k(s,a)
\geq
\max_a Q_h^*(s,a)
=
V_h^*(s).
\]
归纳完成。
\end{proof}

\textbf{Step 3：Regret分解}

由乐观性，
\begin{align*}
\operatorname{Regret}(K)
&=
\sum_{k=1}^K
\left[
V_1^*(s_1^k)-V_1^{\pi^k}(s_1^k)
\right]
\\
&\leq
\sum_{k=1}^K
\left[
V_1^k(s_1^k)-V_1^{\pi^k}(s_1^k)
\right].
\end{align*}

下面分析单个 episode 中
\[
V_1^k(s_1^k)-V_1^{\pi^k}(s_1^k)
\]
的大小。

\begin{lemma}[单episode分解]
在好事件 $\mathcal{E}$ 下，对每个 episode \(k\)，有
\[
V_1^k(s_1^k)-V_1^{\pi^k}(s_1^k)
\leq
\sum_{h=1}^H
\left[
2\operatorname{bonus}_h^k(s_h^k,a_h^k)
+
\xi_h^k
\right],
\]
其中 \(\xi_h^k\) 是鞅差序列。
\end{lemma}

\begin{proof}[引理证明]
定义
\[
\delta_h^k(s)
=
V_h^k(s)-V_h^{\pi^k}(s).
\]
在 episode \(k\) 的真实轨迹上，记
\[
\delta_h^k
=
\delta_h^k(s_h^k).
\]

由于策略 \(\pi^k\) 在第 \(h\) 步选择
\[
a_h^k=\arg\max_a Q_h^k(s_h^k,a),
\]
有
\begin{align*}
\delta_h^k
&=
V_h^k(s_h^k)-V_h^{\pi^k}(s_h^k)
\\
&=
Q_h^k(s_h^k,a_h^k)
-
\left[
R(s_h^k,a_h^k)
+
\sum_{s'}P(s'|s_h^k,a_h^k)
V_{h+1}^{\pi^k}(s')
\right]
\\
&=
\sum_{s'}
\hat P^k(s'|s_h^k,a_h^k)V_{h+1}^k(s')
-
\sum_{s'}
P(s'|s_h^k,a_h^k)V_{h+1}^{\pi^k}(s')
+
\operatorname{bonus}_h^k(s_h^k,a_h^k).
\end{align*}

将第一项加减
\[
\sum_{s'}P(s'|s_h^k,a_h^k)V_{h+1}^k(s')
\]
得到
\begin{align*}
\delta_h^k
&=
\sum_{s'}
(\hat P^k-P)(s'|s_h^k,a_h^k)V_{h+1}^k(s')
\\
&\quad
+
\sum_{s'}
P(s'|s_h^k,a_h^k)
\left[
V_{h+1}^k(s')-V_{h+1}^{\pi^k}(s')
\right]
+
\operatorname{bonus}_h^k(s_h^k,a_h^k)
\\
&\leq
2\operatorname{bonus}_h^k(s_h^k,a_h^k)
+
\mathbb{E}
\left[
\delta_{h+1}^k(s_{h+1}^k)
\mid
s_h^k,a_h^k
\right],
\end{align*}
其中最后一步使用了好事件。

定义鞅差
\[
\xi_h^k
=
\delta_{h+1}^k(s_{h+1}^k)
-
\mathbb{E}
\left[
\delta_{h+1}^k(s_{h+1}^k)
\mid
s_h^k,a_h^k
\right].
\]
于是
\[
\mathbb{E}
\left[
\delta_{h+1}^k(s_{h+1}^k)
\mid
s_h^k,a_h^k
\right]
=
\delta_{h+1}^k(s_{h+1}^k)-\xi_h^k.
\]
因此
\[
\delta_h^k
\leq
2\operatorname{bonus}_h^k(s_h^k,a_h^k)
+
\delta_{h+1}^k(s_{h+1}^k)
-
\xi_h^k.
\]

将上式从 \(h=1\) 递推到 \(H\)，并利用
\[
\delta_{H+1}^k=0,
\]
得到
\[
\delta_1^k
\leq
\sum_{h=1}^H
\left[
2\operatorname{bonus}_h^k(s_h^k,a_h^k)
-
\xi_h^k
\right].
\]
由于 \(-\xi_h^k\) 仍然是鞅差项，为了记号简洁，重新记作 \(\xi_h^k\)，得到
\[
\delta_1^k
\leq
\sum_{h=1}^H
\left[
2\operatorname{bonus}_h^k(s_h^k,a_h^k)
+
\xi_h^k
\right].
\]
\end{proof}

\textbf{Step 4：累积bonus的界}

由 Step 3，
\[
\operatorname{Regret}(K)
\leq
2\sum_{k=1}^K\sum_{h=1}^H
\operatorname{bonus}_h^k(s_h^k,a_h^k)
+
\sum_{k=1}^K\sum_{h=1}^H
\xi_h^k.
\]

首先处理鞅差项。由于
\[
|\xi_h^k|\leq H,
\]
由 Azuma-Hoeffding 不等式，以高概率有
\[
\left|
\sum_{k=1}^K\sum_{h=1}^H
\xi_h^k
\right|
\leq
O\left(H\sqrt{HK\iota}\right).
\]

下面处理 bonus 项。由定义，
\[
\sum_{k=1}^K\sum_{h=1}^H
\operatorname{bonus}_h^k(s_h^k,a_h^k)
=
cH\sqrt{\iota}
\sum_{k=1}^K\sum_{h=1}^H
\frac{1}{
\sqrt{N^k(s_h^k,a_h^k)\vee 1}
}.
\]

对每个固定的 \((s,a)\)，设它在所有 \(K\) 个 episode 中总共被访问了
\[
n(s,a)
\]
次，并把这些访问时刻依次记为
\[
t_1,t_2,\ldots,t_{n(s,a)}.
\]
在第 \(i\) 次访问时，有
\[
N^k(s,a)=i-1.
\]
因此
\[
\sum_{i=1}^{n(s,a)}
\frac{1}{\sqrt{(i-1)\vee 1}}
\leq
1+\sum_{i=1}^{n(s,a)-1}\frac{1}{\sqrt{i}}
\leq
2\sqrt{n(s,a)}.
\]

于是
\begin{align*}
\sum_{k=1}^K\sum_{h=1}^H
\frac{1}{
\sqrt{N^k(s_h^k,a_h^k)\vee 1}
}
&\leq
\sum_{s,a}2\sqrt{n(s,a)}
\\
&\leq
2\sqrt{SA}
\sqrt{\sum_{s,a}n(s,a)}
\\
&=
2\sqrt{SA}
\sqrt{HK}
\\
&=
2\sqrt{SAHK}.
\end{align*}
其中第二步使用 Cauchy-Schwarz，不再对 \(h\) 分桶；并且
\[
\sum_{s,a}n(s,a)=HK.
\]

因此
\[
\sum_{k=1}^K\sum_{h=1}^H
\operatorname{bonus}_h^k(s_h^k,a_h^k)
\leq
O\left(
H\sqrt{SAHK\iota}
\right)
=
O\left(
\sqrt{H^3SAK\iota}
\right).
\]

综上，
\begin{align*}
\operatorname{Regret}(K)
&\leq
O\left(
\sqrt{H^3SAK\iota}
\right)
+
O\left(
H\sqrt{HK\iota}
\right).
\end{align*}
当 \(SA\geq 1\) 时，第一项主导，因此
\[
\operatorname{Regret}(K)
\leq
O\left(
\sqrt{H^3SAK\iota}
\right).
\]
证明完成。
\end{proof}

\begin{calculation}[UCB-VI Regret数值验证]
$|S|=100$，$|A|=4$，$H=20$，$K=10000$，$\delta=0.05$：

$\iota = \ln(100 \times 4 \times 20 \times 10000 / 0.05) = \ln(1.6 \times 10^9) \approx 21$

$\text{Regret} \leq c\sqrt{20^3 \times 100 \times 4 \times 10000 \times 21}$

$= c\sqrt{8000 \times 400 \times 10000 \times 21} = c\sqrt{6.72 \times 10^{11}} \approx 820000 \cdot c$

取$c = 1$（实际常数更小），平均每episode Regret $\approx 82$。

实际中Regret通常比理论界小5-10倍。
\end{calculation}

\subsection{Regret下界}

\begin{theorem}[信息论下界]
存在MDP使得任何算法的期望Regret至少：
\[
\mathbb{E}[\text{Regret}(K)] \geq \Omega(\sqrt{H S A K})
\]
\end{theorem}

\begin{proof}
\textbf{构造困难实例}：

考虑"链式MDP"：
\begin{itemize}
    \item $S$个状态排成一条链：$s_1 \to s_2 \to \cdots \to s_S$
    \item 每个状态有$A$个动作
    \item 只有一个"正确"动作能前进，其他动作原地不动
    \item 正确动作是随机选的（算法不知道）
    \item 到达$s_S$获得奖励1，否则为0
\end{itemize}

\textbf{分析}：

要到达$s_S$，必须在\textbf{每个}状态$s_i$找到正确动作。

在状态$s_i$，需要$\Omega(A)$次尝试才能找到正确动作（每个动作概率$1/A$）。

总共需要$\Omega(SA)$次探索才能学会完整策略。

在探索期间，每个episode的期望Regret是$\Omega(1)$（因为很可能到不了$s_S$）。

设探索需要$T$个episode，则Regret $\geq \Omega(T)$。

但我们也需要$T \cdot H$次转移来完成探索。每次转移最多给一个$(s,a)$提供信息。

要区分每个$(s,a)$是否"正确"，由上周的信息论下界，需要$\Omega(1)$次访问。

总共需要$\Omega(SA)$次访问，即$T \cdot H \geq \Omega(SA)$，所以$T \geq \Omega(SA/H)$。

但这还不够紧。更精细的分析：

\textbf{关键洞察}：在$K$个episode中，每个$(s,a,h)$的访问次数约为$K/(SA)$（平均分配）。

用上周Lai-Robbins风格的论证：要区分两个相差$\Delta$的分布，需要$\Omega(1/\Delta^2)$样本。

对手可以选择$\Delta = \Theta(1/\sqrt{K/(SA)}) = \Theta(\sqrt{SA/K})$。

此时单个状态的Regret贡献$\approx \Delta \cdot K/(SA) \cdot (SA) = \Delta K = \Theta(\sqrt{SAK})$。

乘以$H$步的误差传播：$\text{Regret} \geq \Omega(\sqrt{HSAK})$。
\end{proof}

\begin{keypoint}[上界与下界的gap]
\begin{center}
\begin{tabular}{l|c}
\toprule
界 & Regret \\
\midrule
UCB-VI上界 & $\tilde{O}(\sqrt{H^3 SAK})$ \\
下界 & $\Omega(\sqrt{HSAK})$ \\
Gap & $H$的指数：$1.5$ vs $0.5$ \\
\bottomrule
\end{tabular}
\end{center}

Gap在$H$上。近年有改进算法达到$\tilde{O}(\sqrt{H^2 SAK})$，但$\sqrt{H}$的gap是否能消除仍是开放问题。
\end{keypoint}


%=============================================================================
% 新增章节：时序差分学习——从prediction到control
% 插入位置：§UCB-VI 之后、§Q-learning收敛性 之前
% 配套修改：原§Q-learning可瘦身，"有模型vs无模型"小节可删（已在此节阐明）
%=============================================================================

\section{时序差分学习：从prediction到control}
%=============================================================================

\begin{plainwords}
回顾迄今为止处理MDP的三条路：
\begin{itemize}
\item \textbf{DP（值迭代/策略迭代）}：要$P, R$的解析形式，每步算期望
\item \textbf{Monte Carlo}：model-free，但要\textbf{等整段episode}才能更新，方差大
\item \textbf{UCB-VI}：从样本估出$\hat P, \hat R$再planning（model-based，乐观探索）
\end{itemize}

本节给出第四条路：\textbf{时序差分（Temporal Difference, TD）学习}。它是model-free（像MC）但用\textbf{bootstrap}绕过等待（像DP）：用一步样本$+$当前估计代替整段$G_t$。

\textbf{TD = bandit式样本 + DP式bootstrap}，恰好填上MC和DP之间的空位。

Q-learning和SARSA都不是凭空出现的算法——它们是TD框架在control问题上的两个具体实例，差别只在target里$a'$怎么选（按$\pi$ vs.\ 取$\max$）。
\end{plainwords}

\subsection{Prediction问题：给定策略，估计价值}

\begin{plainwords}[问题设定]
固定一个策略$\pi$。我们\textbf{不知道}$P, R$，但可以让$\pi$在环境中跑，观察序列：
\[
s_0 \xrightarrow{a_0\sim\pi(\cdot|s_0)} r_0, s_1 \xrightarrow{a_1\sim\pi(\cdot|s_1)} r_1, s_2 \xrightarrow{} \cdots
\]
目标：估计$V^\pi(s)$。

注意这\textbf{不是}求最优策略——是给定$\pi$估它的值。术语上区分：
\begin{itemize}
\item \textbf{Prediction（评估）}：给$\pi$估$V^\pi$或$Q^\pi$
\item \textbf{Control（控制）}：找最优$\pi^*$
\end{itemize}
本节先讲prediction（更简单、更基础），再过渡到control。
\end{plainwords}

\begin{keypoint}[估计$V^\pi$的三条路]
回忆Bellman期望方程：
\[
V^\pi(s) = \mathbb{E}_{a\sim\pi}\Big[R(s,a) + \gamma\, \mathbb{E}_{s'\sim P(\cdot|s,a)}[V^\pi(s')]\Big]
\]

右边有两层期望（动作 + 下一状态），还有$V^\pi$自身的递归。处理这三件事的方式不同，得到三种算法：

\begin{center}
\begin{tabular}{l|c|c|l}
\toprule
方法 & 要$P,R$？ & 用样本？ & 怎么处理$V^\pi$的递归 \\
\midrule
DP（策略评估） & 是 & 否 & 用前一轮估计$V_k$（bootstrap） \\
Monte Carlo & 否 & 是（整段） & \textbf{不递归}，直接用真实$G_t$ \\
TD(0) & 否 & 是（一步） & 用当前估计$V(s')$（bootstrap） \\
\bottomrule
\end{tabular}
\end{center}

\textbf{两个轴}：是否用样本（vs.\ 解析期望）、是否bootstrap（vs.\ 用真实回报）。DP和TD都bootstrap，DP用期望、TD用样本；MC不bootstrap但用样本。\textbf{TD = bandit式样本 + DP式bootstrap}。
\end{keypoint}

\subsection{TD(0)：用 bootstrap 替换 $G_t$}

\begin{plainwords}[回顾MC的痛点]
上一章MC的核心痛点是\textbf{要等}：从$s_t$出发，必须把整段轨迹跑完才能算$G_t$，才能更新$V(s_t)$。

TD的诞生动机就是：能不能不等？\textbf{一步}就更新？
\end{plainwords}

TD的核心idea极其简单：

\begin{intuition}[Return的递归关系]
$G_t = r_t + \gamma G_{t+1}$

如果手头已经有$V(s_{t+1})$作为$\mathbb{E}[G_{t+1}|s_{t+1}]$的估计，就可以\textbf{用它代替未知的$G_{t+1}$}：
\[
G_t \approx r_t + \gamma V(s_{t+1})
\]

一步就能更新，不用等到episode结束。这个"用自己的估计替换真值"的动作，就叫\textbf{bootstrap}。
\end{intuition}

\begin{definition}[TD(0)算法——prediction版]
观察到$(s_t, a_t, r_t, s_{t+1})$后：
\[
\boxed{V(s_t) \leftarrow V(s_t) + \alpha\Big[\underbrace{r_t + \gamma V(s_{t+1})}_{\text{TD target}} - V(s_t)\Big]}
\]

定义\textbf{TD误差}：$\delta_t := r_t + \gamma V(s_{t+1}) - V(s_t)$。

更新简写为$V(s_t) \leftarrow V(s_t) + \alpha\,\delta_t$。
\end{definition}

\begin{keypoint}[Bootstrapping的正式定义]
一个学习算法\textbf{bootstrap}，当且仅当：它的更新目标含\textbf{当前估计的价值函数}$V$或$Q$。

\begin{itemize}
\item DP（策略评估/值迭代）bootstrap：$V_{k+1}(s) = \mathbb{E}[r + \gamma V_k(s')]$，target里有$V_k$
\item TD(0) bootstrap：target = $r + \gamma V(s')$，含$V$
\item MC \textbf{不} bootstrap：target = $G_t$，全是真实奖励，不含$V$
\end{itemize}

\textbf{Bootstrap是一把双刃剑}：
\begin{itemize}
\item \textbf{好处}：一步即可更新（不等episode）；方差降低（不叠加整段噪声）
\item \textbf{代价}：引入偏差（$V$估计有错，target也错）；和\textbf{函数逼近 + off-policy}组合时会发散（§Deadly Triad）
\end{itemize}

记住这个定义——后面讨论Deadly Triad、advantage estimation、GAE，"bootstrap"都是核心词。
\end{keypoint}

\begin{example}[MC vs TD：一个具体轨迹]
轨迹：$s_0 \to (r_0=1) \to s_1 \to (r_1=2) \to s_2 \to (r_2=4) \to$ 终止。$\gamma=1$。

当前估计：$V(s_0)=10, V(s_1)=5, V(s_2)=3$。$\alpha=0.1$。

\textbf{MC对$s_0$的更新}：
\begin{itemize}
\item $G_0 = 1+2+4 = 7$（真实回报）
\item Target = 7，误差 = $7-10 = -3$
\item $V(s_0) \leftarrow 10 + 0.1\times(-3) = 9.7$
\end{itemize}

\textbf{TD(0)对$s_0$的更新}（在看到$s_1$后立刻发生，不等episode结束）：
\begin{itemize}
\item Target = $r_0 + \gamma V(s_1) = 1 + 5 = 6$
\item TD误差 $\delta_0 = 6 - 10 = -4$
\item $V(s_0) \leftarrow 10 + 0.1\times(-4) = 9.6$
\end{itemize}

\textbf{两个target不同}：MC用$7$（真实但带$r_1, r_2$的方差），TD用$6$（用$V(s_1)=5$代替了未来——如果这个$5$是准的、target偏差就小，但只承担一步的奖励噪声）。
\end{example}

\begin{calculation}[为什么TD降方差]
\[
\text{Var}(\text{MC target}) = \text{Var}\left(\sum_{k=0}^\infty \gamma^k r_{t+k}\right)
\]
\[
\text{Var}(\text{TD target}) = \text{Var}\big(r_t + \gamma V(s_{t+1})\big)
\]

只要$V$是个相对稳定的估计，TD target只承担\textbf{一步奖励}的随机性$+$ 一个不太抖的$V$值，方差远小于MC target那个无穷和。

代价：若$V(s_{t+1})$有偏差$b$，TD target有偏差$\gamma b$。
\textbf{Bias-variance tradeoff——TD($\lambda$)正是要在这个权衡上插值（见后面§TD($\lambda$)）。}
\end{calculation}

\begin{theorem}[TD(0)收敛性 - Sutton 1988; Dayan 1992]
表格情形下，若：
\begin{enumerate}
\item 所有状态$s$被无限次访问
\item 学习率满足$\sum_t \alpha_t(s) = \infty$，$\sum_t \alpha_t(s)^2 < \infty$
\end{enumerate}
则TD(0)以概率1收敛到$V^\pi$。
\end{theorem}

\begin{proofidea}
TD(0)期望更新方向是
\[
\mathbb{E}[\delta_t \mid s_t = s] = (T^\pi V)(s) - V(s)
\]
其中$T^\pi$是Bellman期望算子——已在§3证明它是$\gamma$-收缩的不动点$V^\pi$。

所以TD(0)是$V_{k+1} = T^\pi V_k$这个DP迭代的\textbf{随机逼近版本}：用样本$(s_t, r_t, s_{t+1})$估计$T^\pi V$。

收敛性证明结构和Q-learning相同（§后面）：收缩算子 + 鞅差噪声 + Robbins-Monro步长 → w.p.1收敛。把Q-learning证明里的$T^*$换成$T^\pi$、$\max_{a'}$换成$\mathbb{E}_{a'\sim\pi}$即可。
\end{proofidea}

\begin{plainwords}[三种prediction方法的统一视图]
\begin{center}
\begin{tabular}{l|l|l|l}
\toprule
& Target & 用模型 & 用样本 \\
\midrule
DP（策略评估） & $\sum_a \pi(a|s)[R + \gamma\sum_{s'} P V_k(s')]$ & 是 & 否 \\
MC评估 & $G_t$（完整回报） & 否 & 是 \\
TD(0) & $r_t + \gamma V(s_{t+1})$ & 否 & 是（一步） \\
\bottomrule
\end{tabular}
\end{center}

\textbf{TD = DP的递归 + MC的样本}。把DP里的"求期望"换成"采样"，把MC里的"完整回报"换成"一步样本+bootstrap"。
\end{plainwords}

\subsection{从prediction到control：为什么必须学$Q$不学$V$}

Prediction解决了"给定$\pi$估值"。要找最优策略呢？

回忆策略迭代的improvement步骤：
\[
\pi'(s) = \arg\max_a Q^\pi(s, a)
\]

\begin{plainwords}[Model-free control的根本约束]
要在不知道$P, R$的情况下做策略改进，必须能比较\textbf{不同动作}的好坏：
\[
\pi'(s) = \arg\max_a Q(s, a)
\]
如果手头只有$V$，想算$Q$需要：
\[
Q^\pi(s, a) = R(s, a) + \gamma \sum_{s'} P(s'|s, a) V^\pi(s')
\]
——这\textbf{需要$P$和$R$}！

所以：\textbf{model-free control必须直接学$Q$，不学$V$}。这是从prediction（学$V$就够）到control（必须学$Q$）的根本驱动力。
\end{plainwords}

把TD(0)从$V$推广到$Q$。需要的样本变为五元组$(s_t, a_t, r_t, s_{t+1}, a_{t+1})$（多了"下一步动作"）：

\begin{definition}[TD(0) for $Q$]
\[
Q(s_t, a_t) \leftarrow Q(s_t, a_t) + \alpha\big[r_t + \gamma Q(s_{t+1}, a_{t+1}) - Q(s_t, a_t)\big]
\]
\end{definition}

\textbf{这里$a_{t+1}$从哪来？}这个看似平凡的选择，决定了你得到的是on-policy还是off-policy算法。两个不同的答案对应SARSA和Q-learning。

\subsection{SARSA：on-policy TD control}

\begin{definition}[SARSA]
名字来自更新所用的五元组：$(s_t, a_t, r_t, s_{t+1}, a_{t+1})$。
\begin{algorithm}[H]
\caption{SARSA}
\begin{algorithmic}[1]
\State 初始化$Q(s,a)$任意
\State 在$s_0$用$\epsilon$-greedy（基于当前$Q$）选$a_0$
\For{每步$t$}
    \State 执行$a_t$，观察$r_t, s_{t+1}$
    \State 在$s_{t+1}$用$\epsilon$-greedy选$a_{t+1}$
    \State $Q(s_t, a_t) \leftarrow Q(s_t, a_t) + \alpha[r_t + \gamma Q(s_{t+1}, a_{t+1}) - Q(s_t, a_t)]$
    \State $t \leftarrow t+1$
\EndFor
\end{algorithmic}
\end{algorithm}

关键：$a_{t+1}$是\textbf{即将真正执行的下一步动作}，按当前behavior policy（$\epsilon$-greedy）采样。
\end{definition}

\begin{keypoint}[SARSA是on-policy的]
SARSA的TD target里的$a_{t+1}\sim \pi_{\text{behavior}}$——\textbf{产生数据的策略}和\textbf{被评估的策略}是同一个。

所以SARSA学到的$Q$，是当前$\epsilon$-greedy策略本身的$Q^\pi$（其中$\pi$ = $\epsilon$-greedy w.r.t.\ $Q$）。

要让$Q \to Q^*$，必须让$\epsilon \to 0$（\textbf{GLIE}: Greedy in the Limit with Infinite Exploration）：
\begin{itemize}
\item 长期$\epsilon \to 0$ ⟹ behavior policy → greedy
\item 但$\epsilon$下降不能太快，否则探索不充分（条件1失败）
\item 常用：$\epsilon_t = 1/t$
\end{itemize}
\end{keypoint}

\subsection{Q-learning：off-policy TD control}

现在做一个微小但关键的修改：把SARSA target里的$Q(s_{t+1}, a_{t+1})$ 换成$\max_{a'} Q(s_{t+1}, a')$。

\begin{definition}[Q-learning]
\[
Q(s_t, a_t) \leftarrow Q(s_t, a_t) + \alpha\big[r_t + \gamma \max_{a'} Q(s_{t+1}, a') - Q(s_t, a_t)\big]
\]
\end{definition}

\begin{magicbox}[$\max$的语义：off-policy的诞生]
SARSA的target $Q(s_{t+1}, a_{t+1})$对应\textbf{当前$\epsilon$-greedy策略}的下一步价值——评估的就是行为策略本身。

Q-learning的target $\max_{a'} Q(s_{t+1}, a')$对应一个\textbf{不同的策略}：
\[
\pi_{\text{greedy}}(s) = \arg\max_a Q(s, a) \quad (\text{无探索的纯贪婪})
\]
但智能体实际\textbf{行动}时仍用$\epsilon$-greedy（必须有探索才能更新所有$(s,a)$）。

\textbf{这就是off-policy的根源}：
\begin{itemize}
\item Behavior policy（产生数据）：$\epsilon$-greedy
\item Target policy（被学习/评估）：greedy
\item 两者不同 ⟹ off-policy
\end{itemize}

\textbf{为什么这样做有意义}：Q-learning学到的是$Q^*$（最优策略的Q），与你用什么探索策略\textbf{无关}（前提是所有$(s,a)$都能被访问到）。SARSA要学$Q^*$得让$\epsilon\to 0$；Q-learning直接学$Q^*$，$\epsilon$只影响样本质量、不影响目标。这是Q-learning相对SARSA最大的优势。代价：函数逼近下会触发Deadly Triad（§后面）。
\end{magicbox}

\begin{keypoint}[SARSA vs Q-learning——同源异果]
\begin{center}
\begin{tabular}{l|c|c}
\toprule
& \textbf{SARSA} & \textbf{Q-learning} \\
\midrule
TD target & $r + \gamma Q(s', a')$ & $r + \gamma \max_{a'} Q(s', a')$ \\
$a'$ 怎么选 & 从$\pi_{\text{behavior}}$采样 & 取$\arg\max$ \\
On/off-policy & on-policy & off-policy \\
评估的策略 & behavior policy（$\epsilon$-greedy） & greedy policy \\
学到什么 & $Q^{\pi_{\text{behavior}}}$ & $Q^*$ \\
何时$\to Q^*$ & 需要GLIE（$\epsilon\to 0$） & 直接收敛（只要访问充分） \\
表格收敛性 & w.p.1（GLIE下） & w.p.1（Watkins 1992） \\
\bottomrule
\end{tabular}
\end{center}

两个算法\textbf{只差一个$\max$}，但语义完全不同——一个是"评估自己（含探索）"，一个是"评估假想的greedy自己"。
\end{keypoint}

\begin{example}[悬崖行走：on-policy vs off-policy的实际差别]
经典对比（Sutton \& Barto, Example 6.6）：智能体在网格世界，要从S走到G。地图底部一长条是"悬崖"，掉下去$-100$奖励并送回起点；其他每步$-1$。

\textbf{Q-learning}学到的策略：贴着悬崖边走最短路径（这是$\pi^*$）。但训练过程中智能体按$\epsilon$-greedy行动，所以经常被$\epsilon$的扰动推下悬崖，\textbf{训练期累积奖励反而差}。

\textbf{SARSA}学到的策略：绕远路、远离悬崖。因为SARSA评估的就是$\epsilon$-greedy策略本身——它"知道"自己以后还会随机出错，所以选一条容错率高的路。\textbf{训练期累积奖励高，但学到的策略不是$\pi^*$。}

\textbf{道理}：\textbf{Q-learning假装自己最终会完美，SARSA承认自己永远会犯小错}。两者各有适用场景：要部署时关探索，Q-learning好；要训练中也少出大事故（机器人、医疗），SARSA好。
\end{example}

\begin{plainwords}[Tabular RL算法的二维分类表]
现在我们终于可以把表格RL所有算法放在一张图里：

\begin{center}
\begin{tabular}{l|l|l}
\toprule
& \textbf{Prediction（评估$\pi$）} & \textbf{Control（找$\pi^*$）} \\
\midrule
\textbf{DP}（要模型） & 策略评估 & 值迭代 / 策略迭代 \\
\textbf{MC}（无模型、不bootstrap） & MC评估 & MC control + $\epsilon$-greedy \\
\textbf{TD on-policy} & TD(0) on $V$ & \textbf{SARSA} \\
\textbf{TD off-policy} & 重要性采样TD & \textbf{Q-learning} \\
\bottomrule
\end{tabular}
\end{center}

\textbf{记忆要点}：
\begin{itemize}
\item 横轴：估值 vs 找最优
\item 纵轴：怎么处理期望（解析 / 整段样本 / 一步样本+bootstrap）和on/off-policy
\item Q-learning $=$ TD $+$ control $+$ off-policy（$\max$）
\item SARSA $=$ TD $+$ control $+$ on-policy（按$\pi$采样$a'$）
\end{itemize}

后面要讲的所有东西都长在这张表上：
\begin{itemize}
\item TD($\lambda$)：在TD(0)（一步）和MC（整段）之间用$\lambda$插值
\item Linear TD/Q：把表格换成$\theta^\top\phi(s)$
\item Deadly Triad：FA + bootstrap + off-policy组合发散
\item Advantage function / policy gradient（后续课）：TD误差$\delta_t$恰是advantage $A^\pi(s_t,a_t)$的无偏估计
\item Deep RL的DQN：Q-learning + 神经网络 + replay buffer + target network
\end{itemize}
\end{plainwords}

\begin{remark}[术语对照——避免课后混淆]
不同教材、论文对同一概念叫法不一，列在这里：
\begin{itemize}
\item \textbf{Prediction} = policy evaluation = 策略评估
\item \textbf{Control} = optimization = 找最优策略
\item \textbf{Bootstrap} = 用估计更新估计（DP和TD都bootstrap，MC不）
\item \textbf{On-policy}：target policy = behavior policy
\item \textbf{Off-policy}：target policy $\neq$ behavior policy
\item \textbf{TD error} $\delta_t = r_t + \gamma V(s_{t+1}) - V(s_t)$（或Q版本）
\item \textbf{TD target} $= r_t + \gamma V(s_{t+1})$（TD误差里$V(s_t)$之外的部分）
\end{itemize}
\end{remark}

%=============================================================================
% 下一节：§Q-learning收敛性（原§8，可瘦身——
% 删去"有模型vs无模型"的1.5页（已在本节阐明）、"试错学习"那段比喻、
% 只保留：算法伪码（已在本节有）、Watkins-Jaakkola定理 + 证明）
%=============================================================================

%=============================================================================
\section{核心概念与工程技巧：on/off-policy、经验回放、目标网络}
%=============================================================================

\begin{plainwords}
到这里我们已经见过一个不小的算法动物园：值迭代、策略迭代、MC、UCB-VI、TD(0)、SARSA、Q-learning。它们看起来五花八门，但其实坐在同一个低维空间里——只要找对了坐标轴。

本节做两件事：

\textbf{(1) 概念坐标}。给出区分RL算法的三个最重要维度：
\begin{itemize}
\item \textbf{model-based vs model-free}（要不要$P, R$）——前几节的主线
\item \textbf{on-policy vs off-policy}（执行的策略 = 评估的策略 吗）
\item \textbf{online vs offline}（边交互边学 vs 只用历史数据）
\end{itemize}

\textbf{(2) 三件工程套件}。介绍把上面算法做"实战"的三个核心技巧：
\begin{itemize}
\item \textbf{经验回放（experience replay）}——让 online 算法重复消化样本
\item \textbf{目标网络（target network）}——让bootstrap稳定
\item \textbf{重要性采样（importance sampling）}——off-policy的统计桥梁
\end{itemize}

外加一个常见陷阱：\textbf{Q-learning的过估计偏差}与\textbf{Double Q-learning}。

这些是后面看任何deep RL文献都绕不开的词汇。
\end{plainwords}

\subsection{On-policy vs off-policy：精确定义}

\begin{definition}[Behavior policy 与 target policy]
任何RL算法在某个时刻都涉及两个策略：
\begin{itemize}
\item \textbf{Behavior policy} $\pi_b$：\textbf{产生数据}的策略（"我实际怎么走、观察到了什么"）
\item \textbf{Target policy} $\pi_t$：\textbf{被学习/评估}的策略（"我想知道哪个策略的价值/我想优化哪个策略"）
\end{itemize}

\textbf{On-policy}：$\pi_b = \pi_t$（评估和执行是同一个策略）

\textbf{Off-policy}：$\pi_b \neq \pi_t$（用一个策略的数据学习另一个策略）
\end{definition}

\begin{example}[把已学算法摆进去]
\begin{center}
\begin{tabular}{l|c|c|l}
\toprule
算法 & $\pi_b$ & $\pi_t$ & 是哪一种 \\
\midrule
MC评估 & $\pi$ & $\pi$ & on-policy \\
MC control & $\epsilon$-greedy & $\epsilon$-greedy & on-policy \\
SARSA & $\epsilon$-greedy & $\epsilon$-greedy & on-policy \\
Q-learning & $\epsilon$-greedy & greedy & \textbf{off-policy} \\
UCB-VI & 乐观greedy & 乐观greedy & on-policy（执行=评估）\\
\bottomrule
\end{tabular}
\end{center}

\textbf{Q-learning成为off-policy的唯一原因}：target里的$\max_{a'} Q(s', a')$对应的是greedy策略（$\pi_t$），与实际行动的$\epsilon$-greedy（$\pi_b$）不同。改成target里按$\pi_b$采样$a'$，立刻变回on-policy（即SARSA）。
\end{example}

\begin{keypoint}[Off-policy的好处]
\begin{enumerate}
\item \textbf{探索与利用解耦}：用充分探索的$\pi_b$产生数据，学到完全greedy的$\pi_t$。最优行为不被探索噪声污染。
\item \textbf{数据复用}：旧数据（不管由哪个旧策略产生）可以反复利用——这是\textbf{经验回放}的前提
\item \textbf{演示/示教学习}：可以从人类专家或别的策略的数据中学习
\item \textbf{多策略并行评估}：用同一份数据同时评估多个候选target policy
\end{enumerate}
\end{keypoint}

\begin{keypoint}[Off-policy的代价]
\begin{enumerate}
\item \textbf{分布不匹配}：$\pi_b$产生的状态/动作分布$d^{\pi_b} \neq d^{\pi_t}$，朴素估计会有偏。修正工具：重要性采样（下面讲）
\item \textbf{Bellman最优算子的max带来过估计}（见后§Double Q-learning）
\item \textbf{和函数逼近 + bootstrap 组合 = Deadly Triad}（见§Linear FA）——deep RL最常踩的坑
\end{enumerate}
\end{keypoint}

\subsection{重要性采样：off-policy的统计桥梁}

\begin{plainwords}
要从$\pi_b$产生的样本估计$\pi_t$下的期望，最干净的统计工具是\textbf{重要性采样（importance sampling, IS）}。

直觉：用比率$\rho = \pi_t/\pi_b$给每个样本加权——被$\pi_b$多观察到的事件权重降低，反之升高。
\end{plainwords}

\begin{theorem}[重要性采样恒等式]
\[
\mathbb{E}_{a \sim \pi_t}[f(a)] = \mathbb{E}_{a \sim \pi_b}\left[\frac{\pi_t(a)}{\pi_b(a)} f(a)\right] = \mathbb{E}_{a \sim \pi_b}[\rho(a) f(a)]
\]
前提（覆盖性条件）：$\pi_t(a) > 0 \Rightarrow \pi_b(a) > 0$。
\end{theorem}

\begin{proof}
\[
\mathbb{E}_{a \sim \pi_t}[f(a)] = \sum_a \pi_t(a) f(a) = \sum_a \pi_b(a) \cdot \frac{\pi_t(a)}{\pi_b(a)} f(a) = \mathbb{E}_{a \sim \pi_b}[\rho(a) f(a)]
\]
两步代数。
\end{proof}

\begin{calculation}[多步IS与方差爆炸]
对$T$步轨迹做off-policy评估，需要乘\textbf{累积似然比}：
\[
\rho_{0:T-1} = \prod_{t=0}^{T-1} \frac{\pi_t(a_t|s_t)}{\pi_b(a_t|s_t)}
\]
这个连乘让方差呈\textbf{指数爆炸}。若每步$\rho_t \sim$均值1、方差$\sigma^2$，则$\text{Var}(\rho_{0:T-1}) \sim (1+\sigma^2)^T - 1$——这就是为什么纯IS只在短horizon可行。
\end{calculation}

\begin{magicbox}[Q-learning为什么不需要IS]
你也许会问：Q-learning是off-policy（$\pi_b = \epsilon$-greedy, $\pi_t = $ greedy），为什么从来没出现似然比？

\textbf{关键}：Q-learning的target $r + \gamma\max_{a'} Q(s', a')$是\textbf{Bellman optimality算子}对$Q$的一步应用，只用到$(s, a, r, s')$这一步转移——\textbf{和产生这个转移的策略无关}（环境dynamics $P, R$不依赖策略）。

所以Q-learning评估的不是某个target policy $\pi_t$的"长期价值"$Q^{\pi_t}$，而是\textbf{最优算子$T^*$的不动点}$Q^*$。这绕开了"用$\pi_b$的轨迹估$\pi_t$的return"那种需要IS校正的设定。

代价：换成了max带来的偏差——见下面的过估计。

这是Q-learning设计上的精妙之处：用\textbf{算子}（关于环境，与策略无关）代替\textbf{轨迹}（依赖策略），完美绕过IS的方差爆炸。
\end{magicbox}

\subsection{Online vs offline：另一个独立维度}

\begin{definition}[Online vs offline]
\textbf{Online}：算法在与环境的\textbf{持续交互}中学习。每步执行动作 → 看反馈 → 更新 → 继续执行。

\textbf{Offline}（也叫batch RL）：只有一个\textbf{预先收集好}的数据集$\mathcal D = \{(s_i, a_i, r_i, s'_i)\}_{i=1}^N$。算法不能再与环境交互，全部学习在$\mathcal D$上完成。
\end{definition}

\begin{keypoint}[Off-policy和offline是不同的轴——最容易混的两个概念]
\begin{center}
\begin{tabular}{l|c|c}
\toprule
& \textbf{Online} & \textbf{Offline} \\
\midrule
\textbf{On-policy} & SARSA, A2C, PPO & \textbf{不存在}（offline下$\pi_b$固定无法 = $\pi_t$） \\
\textbf{Off-policy} & Q-learning, DQN, SAC & CQL, IQL, BCQ, Decision Transformer \\
\bottomrule
\end{tabular}
\end{center}

记住：
\begin{itemize}
\item \textbf{Online $\neq$ on-policy}：Q-learning是online但off-policy
\item \textbf{Offline必然off-policy}：数据已固定，无法再用现在的策略采样
\item \textbf{On-policy + offline = 不存在}（除非定义为supervised learning / 行为克隆）
\end{itemize}
\end{keypoint}

\begin{plainwords}[Offline RL的本质困难]
为什么offline比online难得多？

\textbf{分布偏移（distribution shift）}：你想学的policy $\pi_t$会去$\pi_t$喜欢的状态，但数据集里全是$\pi_b$喜欢的状态。在$\pi_b$从未访问过的$(s, a)$上，$Q$值是\textbf{凭空外推}出来的——而$\max$算子最喜欢挑高估值的那个外推区域，越学越离谱。

Online的Q-learning不怕这个，因为它真去执行$\arg\max$动作、立刻得到环境的反馈纠正。Offline没有这个反馈渠道。

\textbf{对策思路}：让学到的policy"贴近数据分布"——这就是CQL（保守地压低OOD动作的Q值）、IQL（避免显式max操作）、BCQ（限制动作只能从行为分布中选）的共同主线。
\end{plainwords}

\subsection{经验回放：让 online 算法重复消化样本}

\begin{plainwords}
TD/Q-learning的"原生"形式：观察一个$(s,a,r,s')$，更新一次，丢掉，再观察下一个。

两个浪费：
\begin{enumerate}
\item \textbf{样本只用一次}就丢——但收集样本很贵（特别是真实机器人/医疗/工业过程）
\item \textbf{相邻样本强相关}（同一轨迹里前后步状态相似），违反随机逼近定理喜欢的i.i.d.假设，让神经网络训练不稳
\end{enumerate}

\textbf{经验回放（Experience Replay, Lin 1992; Mnih et al.\ 2015）}：把每个观察到的$(s,a,r,s')$存进一个固定大小的buffer，更新时从中随机抽一个mini-batch。一颗样本被反复消化。
\end{plainwords}

\begin{algorithm}[H]
\caption{Q-learning with Experience Replay}
\begin{algorithmic}[1]
\State 初始化$Q$，replay buffer $\mathcal B = \emptyset$（容量$N$，典型$N=10^6$）
\For{每步$t$}
    \State 在$s_t$用$\epsilon$-greedy选$a_t$，执行后观察$r_t, s_{t+1}$
    \State $\mathcal B \leftarrow \mathcal B \cup \{(s_t, a_t, r_t, s_{t+1})\}$（满了FIFO丢老的）
    \State 从$\mathcal B$中均匀采样一个mini-batch $B$（典型$|B|=32\sim 256$）
    \For{$(s, a, r, s') \in B$}
        \State $Q(s,a) \leftarrow Q(s,a) + \alpha[r + \gamma\max_{a'} Q(s', a') - Q(s,a)]$
    \EndFor
\EndFor
\end{algorithmic}
\end{algorithm}

\begin{keypoint}[经验回放为什么有效，以及它的前提]
\textbf{三个好处}：
\begin{enumerate}
\item \textbf{样本复用}：每个$(s,a,r,s')$被使用多次（典型$\sim 8\times$），样本效率近似提升一个数量级
\item \textbf{去相关}：mini-batch里的样本来自不同时刻、不同episode，近似i.i.d.，更接近随机逼近定理的假设
\item \textbf{稳定梯度}：mini-batch平均降低梯度方差，让bootstrap更新更稳
\end{enumerate}

\textbf{关键前提：算法必须是off-policy。}

为什么？buffer里的数据来自\textbf{过去的策略}，按现在的policy可能根本不会做那些动作。on-policy算法（SARSA, PPO）的target要求$a'$按当前$\pi$采样——但你存的是旧$\pi$下的$a'$，分布不对，更新方向有偏。

\textbf{所以经验回放是Q-learning家族的专利，SARSA/PPO没法直接用}。这是on/off-policy差别最实际、最频繁出现的工程后果。
\end{keypoint}

\begin{example}[Prioritized Experience Replay（PER, Schaul et al.\ 2016）]
均匀采样把所有样本一视同仁，但有些样本（TD误差大的）更有学习价值——它们意味着"模型当前预测错得厉害"。

\textbf{PER}：按$|\delta_i|$给样本分配采样概率$p_i \propto |\delta_i|^\alpha$，对$|\delta|$大的样本采得更频繁。需要重要性采样权重$w_i \propto (Np_i)^{-\beta}$校正引入的偏差。

\textbf{效果}：典型场景下样本效率再提升$2\sim 4$倍。代价：需要维护sum-tree数据结构，每次update都要重排。
\end{example}

\subsection{目标网络：bootstrap的稳定化器}

\begin{plainwords}
Q-learning的target：$r + \gamma\max_{a'} Q(s', a')$。注意target里的$Q$\textbf{和被更新的$Q$是同一个}——你在追一个一直在动的靶子。

表格情形下这不要紧（每步只改一个$(s,a)$）。但用神经网络逼近Q时，一次梯度step会同时影响\textbf{所有}$(s, a)$处的$Q$值，target随之大幅波动，bootstrap的偏差被放大、训练经常发散。

\textbf{目标网络（Target Network, Mnih et al.\ 2015）}：维护一个$Q$的\textbf{滞后副本}$Q_{\text{target}}$，在target里用$Q_{\text{target}}$而不是$Q$。
\end{plainwords}

\begin{definition}[带target network的Q-learning更新]
\[
Q(s,a) \leftarrow Q(s,a) + \alpha\Big[r + \gamma\max_{a'} Q_{\text{target}}(s', a') - Q(s,a)\Big]
\]
$Q_{\text{target}}$的两种更新方式：
\begin{itemize}
\item \textbf{硬更新（hard update）}：每$C$步把$Q_{\text{target}} \leftarrow Q$（DQN原版，$C=10^4$）
\item \textbf{软更新（Polyak / soft update）}：每步$Q_{\text{target}} \leftarrow (1-\tau)Q_{\text{target}} + \tau Q$（$\tau \approx 0.005$，常用于DDPG、SAC）
\end{itemize}
\end{definition}

\begin{intuition}[为什么这样稳定]
不动的target = 转化成了\textbf{有监督回归}。在硬更新的两次同步之间，$Q_{\text{target}}$固定不动，所以target $r + \gamma\max Q_{\text{target}}(s')$对每个样本是个\textbf{常数}——拟合$Q$向它逼近就是普通的回归问题，神经网络擅长。每过$C$步换一组新的target，再回归一段。

类比：你在调一个PID控制器，但setpoint一直在跳——调起来当然发散。把setpoint每段时间冻住，反而能调进去。
\end{intuition}

\subsection{Q-learning的过估计偏差与Double Q-learning}

\begin{plainwords}
Q-learning有一个隐蔽的系统性偏差：\textbf{$Q$值会被高估}。原因是max算子和噪声的相互作用。
\end{plainwords}

\begin{theorem}[Jensen不等式：max的偏差]
对随机变量$X_1, \ldots, X_n$，
\[
\mathbb{E}\Big[\max_i X_i\Big] \geq \max_i \mathbb{E}[X_i]
\]
等号仅在所有$X_i$恒等时成立。\textbf{max是凸函数，noise一过它就变正偏。}
\end{theorem}

\begin{plainwords}[Q-learning为什么过估计]
Q-learning的target含$\max_{a'} Q(s', a')$。$Q(s', a')$是带噪声的估计，分解为
\[
Q(s', a') = Q^*(s', a') + \varepsilon_{a'},\qquad \mathbb{E}[\varepsilon_{a'}] = 0
\]
对noise取max：
\[
\max_{a'} Q(s', a') = \max_{a'}\big(Q^*(s', a') + \varepsilon_{a'}\big)
\]
取期望，由Jensen：
\[
\mathbb{E}\big[\max_{a'} Q(s', a')\big] \geq \max_{a'} Q^*(s', a')
\]
即使每个$\varepsilon_{a'}$都零均值，\textbf{被$\max$选中的那个动作的$\varepsilon$倾向于偏正}——因为max偏好挑大噪声的样本。

\textbf{后果}：target被系统性高估 $\to$ $Q$被推高 $\to$ 下次更新更高 $\to$ 一直往上漂，策略对噪声特别敏感，训练有时崩盘。
\end{plainwords}

\begin{definition}[Double Q-learning, van Hasselt 2010]
思路：把\textbf{选动作}和\textbf{评估动作}用两个独立的$Q$估计——一个挑出$\arg\max$，另一个给它估值，noise解耦。

维护两个网络$Q_A, Q_B$。每步随机选一个更新：
\[
Q_A(s, a) \leftarrow Q_A(s,a) + \alpha\Big[r + \gamma\, Q_B\big(s',\,\arg\max_{a'} Q_A(s', a')\big) - Q_A(s,a)\Big]
\]
（再以对称形式更新$Q_B$。）

关键：用$Q_A$\textbf{挑动作}（$\arg\max$），用$Q_B$\textbf{估值}。$Q_B$对$Q_A$挑出来的动作给出\textbf{独立}的估计——max偏差消失，因为max的noise和评估的noise不再相关。

\textbf{Double DQN}（深度版本，Hasselt et al.\ 2016）：用主网络挑动作、target网络估值，几乎免费实现，是DQN系列的标配改进。
\end{definition}

\subsection{探索策略小结}

\begin{plainwords}
我们已经见过几种探索方式，放在一起对比：
\begin{itemize}
\item \textbf{$\epsilon$-greedy}（MC control, SARSA, Q-learning, DQN）：以$\epsilon$概率随机、$1-\epsilon$概率greedy。简单但不智能（在已经熟悉的状态也乱来）
\item \textbf{UCB式bonus}（UCB-VI）：把"不确定性"加到价值估计上，引导探索去不熟悉的$(s,a)$。理论好，但需要计数器
\item \textbf{Boltzmann/softmax}：$\pi(a|s) \propto \exp(Q(s,a)/T)$，温度$T$控制探索强度。比$\epsilon$-greedy更平滑、对Q差异敏感
\item \textbf{Thompson sampling}（贝叶斯RL）：从$Q$的后验分布采样动作。Bandit上经常打败UCB
\item \textbf{内在动机/curiosity}（深度RL）：用状态预测误差等"惊讶度"作内在奖励——处理稀疏外部奖励的利器
\end{itemize}

通用规律：探索越智能（用上更多结构），样本效率越高，但实现复杂度也越高。\textbf{深度RL实际中$\epsilon$-greedy + replay + target net + Double仍是最常见的组合}，简单粗暴管用。
\end{plainwords}

\subsection{Rollout：把"模拟一段轨迹"独立成一个动词}

\begin{plainwords}
"Rollout" 这个词本课前面已经用过很多次但没有正式定义。它就是一个简单的动作：

\textbf{Rollout = 从某个状态 $s$ 出发、按某个策略 $\pi$ 跑一段轨迹（真实或想象的），把过程中的奖励/状态/动作收集起来。}

你已经在每个章节里都做过它：
\begin{itemize}
\item MC评估：从$s$ rollout 到 episode 结束，记录 $G_t$
\item TD(0)：从$s$ rollout \textbf{一步}，剩下用 $V(s')$ bootstrap
\item $n$-step return：从$s$ rollout \textbf{$n$步}，剩下用 $V(s_{t+n})$ bootstrap
\item UCB-VI 内部的值迭代：在估出来的模型 $\hat P, \hat R$ 上做规划——也是一种 rollout，只不过在估的模型里
\end{itemize}

Rollout 是 RL 算法的\textbf{基本动作}，独立于"学什么"和"怎么更新"。明确把它当作概念，能更清楚地看 deep RL 论文里 actor、worker、planner 这些角色。
\end{plainwords}

\begin{keypoint}[Rollout 长度的谱系]
按 rollout 长度区分 RL 算法：
\begin{center}
\begin{tabular}{l|l|l}
\toprule
长度 & Target & 算法 \\
\midrule
1 步 & $r + \gamma V(s')$ & TD(0), SARSA, Q-learning \\
$n$ 步 & $\sum_{k=0}^{n-1}\gamma^k r_{t+k} + \gamma^n V(s_{t+n})$ & $n$-step TD, Rainbow DQN \\
$\lambda$-加权 & $\sum_n (1-\lambda)\lambda^{n-1}G_t^{(n)}$ & TD($\lambda$), GAE \\
整段 & $G_t$ & MC \\
\bottomrule
\end{tabular}
\end{center}

\textbf{Bias-variance 沿 rollout 长度的谱系}：
\begin{itemize}
\item 短 rollout：bootstrap 占比大 → 偏差大（$V$估计不准 target 就不准），方差小（少叠加奖励噪声）
\item 长 rollout：真实奖励占比大 → 偏差小，方差大
\item Rollout 长度本质上是一个连续可调的旋钮——TD(λ) 通过 $\lambda$ 平滑控制
\end{itemize}
\end{keypoint}

\paragraph{Rollout 在 RL 系统里的四种用途}

\textbf{(1) 数据采集（最基本用法）}

Online RL 标准做法：actor 进程在真实环境里 rollout 收集 $(s, a, r, s')$，喂给 learner 进程更新参数。

\begin{itemize}
\item \textbf{同步}（A2C, PPO）：一批 actor 同时 rollout 固定步数，然后 learner 一次性更新
\item \textbf{异步}（A3C, IMPALA, Ape-X, SEED RL）：actor 持续 rollout、learner 持续更新，互不阻塞。需要处理"actor 用的是旧版策略"的 off-policy 校正——IMPALA 的 V-trace 算法是这一线的代表
\end{itemize}

Deep RL 实际跑起来，rollout 经常占 70\%+ 时间——这是为什么 RL infra 设计的核心常常是"如何高效地并行 rollout"。

\textbf{(2) Decision-time planning（MCTS）}

不是把 rollout 用来训练，而是用来\textbf{在当前状态下做决策}。

Monte Carlo Tree Search (MCTS)：在当前状态 $s$ 要决定动作时，模拟很多次（典型 $\sim 10^3-10^5$ 次）"如果我这么走、对手那么走、..."，每次跑到底（或固定深度后用 $V$ bootstrap），统计每个动作的平均回报，选最好的。

\begin{itemize}
\item \textbf{AlphaGo}：用神经网络指导 rollout 中的动作选择（policy network），用神经网络对叶节点估值（value network），UCB-style 选展开节点（UCT）
\item \textbf{AlphaZero / MuZero}：进一步用学到的模型做 rollout（MuZero 甚至不需要环境真实 dynamics，只学一个 latent 模型让 MCTS 在里面跑）
\end{itemize}

MCTS 的精髓：每个动作的价值估计 = 多次 rollout 的样本均值——这其实是 bandit-style sample average 套了一层 tree 结构。

\textbf{(3) Imaginary rollout（model-based RL）}

学到一个模型 $\hat P, \hat R$ 后，可以在\textbf{模型里} rollout 生成"想象的"经验，喂给标准 model-free 算法。

\begin{itemize}
\item \textbf{Dyna-Q}（Sutton 1990）：每次真实交互后，从模型里采样 $k$ 步 rollout 也喂给 Q-learning。最早的 model-based + model-free 混血
\item \textbf{MBPO}（Janner et al.\ 2019）：从 replay buffer 里的真实 $s$ 出发，用学到的模型做\textbf{短 rollout}（避免长期复合误差），把这些 imagined transitions 加进 buffer
\item \textbf{Dreamer}（Hafner et al.\ 2020）：学一个 latent dynamics 模型，把 rollout 完全放在 latent space 做——计算便宜得多
\end{itemize}

核心权衡：模型 rollout 越长，复合误差越大（前面 §Simulation Lemma 分析过：$O(1/(1-\gamma)^2)$ 放大）；越短，越接近 model-free。所以现代 MBRL 一般用\textbf{短 rollout + 真实数据混合}。

\textbf{(4) Rollout algorithm（Tesauro \& Galperin 1996; Bertsekas）}

经典的"做一步策略改进"算法：给定 base policy $\pi_b$，要决定在状态 $s$ 做什么。对每个候选动作 $a$：
\begin{enumerate}
\item 假定先执行 $a$，然后按 $\pi_b$ rollout 到底
\item 多次重复，估计 $\hat Q^{\pi_b}(s, a)$（蒙特卡洛样本均值）
\item 选 $\arg\max_a \hat Q^{\pi_b}(s, a)$
\end{enumerate}

这等价于把策略迭代的"一步改进"做成 online 蒙特卡洛版本。\textbf{保证}：得到的策略至少和 $\pi_b$ 一样好（策略改进定理）。代价：每决策一次都要做一堆 rollout，慢。

历史上 Tesauro 的 TD-Gammon 系列就是用这个把人类水平的 backgammon 程序做到超人。MCTS 可以看作它的"tree-structured"升级版。

\begin{plainwords}[一句话总结四种用法]
\begin{itemize}
\item \textbf{用真实环境 rollout 收集训练数据} → online RL 主流（PPO/DQN/SAC的actor）
\item \textbf{用 rollout 在决策时刻搜索} → MCTS（AlphaGo 系列）
\item \textbf{用模型 rollout 生成合成数据} → model-based RL（Dyna/MBPO/Dreamer）
\item \textbf{用 rollout 做一步策略改进} → Rollout algorithm（经典 Bertsekas 方法）
\end{itemize}

同一个动作（rollout），放在不同位置就是不同算法家族。
\end{plainwords}

\begin{example}[一个状态、四种 rollout 用法的对比]
机械臂抓取任务。当前状态：手在桌面上方某处，眼前有个杯子。

\textbf{(1) 训练数据 rollout}（DQN actor）：执行当前 $\epsilon$-greedy 策略，从这个状态走到 episode 结束，把每步 $(s, a, r, s')$ 塞进 replay buffer。

\textbf{(2) Decision-time rollout}（MCTS）：从这个状态做 $10^4$ 次模拟（在环境模型里），统计每个动作（向左、向右、向下、抓）的平均回报，挑最好的执行。

\textbf{(3) Imaginary rollout}（MBPO）：从 buffer 里采一个真实状态，用学到的 dynamics 模型走 5 步 imaginary trajectory，把这 5 个合成的 transition 也塞进训练 batch。

\textbf{(4) Rollout algorithm}（base policy = "总是向下"）：对每个候选动作 $a$，假定先执行 $a$、然后一直向下，跑 100 次估计 $\hat Q^{\pi_b}(s, a)$，挑最大的 $a$ 执行。每步决策都重新 rollout 一遍。

\textbf{同一个"rollout"动作，承担了完全不同的算法角色}。
\end{example}

\subsection{其他小技巧速览}

\begin{plainwords}
还有几个反复出现的小技巧，不像 replay/target net 那么"必备"，但在 deep RL 论文里几乎处处可见。这一节做个速览，方便后面读论文时不发懵。
\end{plainwords}

\paragraph{(1) $n$-step return：1-step TD 和 MC 之间的可调节插值}

回忆 TD(0) 的 target 是 $r_t + \gamma V(s_{t+1})$（一步真实奖励 + 一步 bootstrap），MC 的 target 是 $G_t = \sum_{k=0}^\infty \gamma^k r_{t+k}$（全部真实奖励）。$n$-step 就在中间插值：
\[
G_t^{(n)} = \underbrace{\sum_{k=0}^{n-1} \gamma^k r_{t+k}}_{\text{$n$步真实奖励}} + \underbrace{\gamma^n V(s_{t+n})}_{\text{$n$步后 bootstrap}}
\]
对应的更新：$V(s_t) \leftarrow V(s_t) + \alpha\big[G_t^{(n)} - V(s_t)\big]$。

\textbf{偏差-方差权衡}：$n$ 大 → 接近 MC，方差大偏差小；$n$ 小 → 接近 TD(0)，偏差大方差小。实践中 $n = 3 \sim 5$ 常常比 1-step 好（Rainbow DQN 用 $n=3$）。更优雅的版本是后面 §TD($\lambda$) 的 $\lambda$-weighted 平均——所有 $n$ 加权混合。

\paragraph{(2) Reward 标准化与 clipping}

不同任务奖励量级差异巨大：CartPole 每步 $+1$，Atari 单关可上千分，机械臂任务奖励可能在 $[-100, 0]$。神经网络（特别是带 Adam 优化器的）对输入/输出尺度敏感，奖励量级直接影响 Q 值量级、TD 误差量级、梯度量级。

\textbf{常见做法}：
\begin{itemize}
\item \textbf{Running normalization}：维护 $r$ 的滑动 mean / std，用 $(r - \mu)/\sigma$ 喂给网络
\item \textbf{Reward clipping}：DQN 在 Atari 上用 $\text{sign}(r) \in \{-1, 0, 1\}$ 暴力 clip——丢失幅度信息，但保留了"好/坏/无"
\item \textbf{Return normalization}（PPO 等）：归一化 $G_t$ 而非单步 $r$
\end{itemize}

\textbf{副作用}：reward clipping 可能扭曲最优策略（10分和100分的区别消失了）。当不同动作的奖励差别在量级上很重要时不要 clip。

\paragraph{(3) Gradient clipping}

TD 误差偶尔会出现 outlier（比如刚开始训练时，或遇到罕见但高奖励的样本），导致梯度突然爆炸、把网络打飞。

\textbf{标准对策}：clip 梯度的 $\ell_2$ 范数：
\[
\nabla \leftarrow \nabla \cdot \min\Big(1,\, \frac{c}{\|\nabla\|_2}\Big)
\]
典型 $c = 0.5 \sim 10$。几乎所有 deep RL 实现都默认开启。

\paragraph{(4) Huber loss：Q-learning 的鲁棒损失}

Q-learning 的标准 MSE loss 在大 TD 误差处会有过大梯度（$\nabla = 2 \delta$），加剧 outlier 的破坏。

\textbf{Huber loss}（DQN 标配）：
\[
L_\kappa(\delta) = \begin{cases} \tfrac{1}{2}\delta^2 & |\delta| \le \kappa \\ \kappa(|\delta| - \tfrac{\kappa}{2}) & |\delta| > \kappa \end{cases}
\]
小误差处像 MSE（光滑、梯度精确），大误差处像 MAE（梯度有界，鲁棒）。典型 $\kappa = 1$。

\paragraph{(5) Dueling 架构：$Q = V + A$ 分头学}

很多状态下"做哪个动作都差不多"——这时只学 $V(s)$ 就够，每个 $Q(s, a)$ 单独学浪费样本。

\textbf{Dueling DQN（Wang et al.\ 2016）}：网络分两个 head，分别输出 $V(s; \theta_v)$ 和 advantage $A(s, a; \theta_a)$，组合：
\[
Q(s, a; \theta) = V(s; \theta_v) + \Big(A(s, a; \theta_a) - \frac{1}{|\mathcal{A}|}\sum_{a'} A(s, a'; \theta_a)\Big)
\]
减均值是为了 identifiability（否则 $V$ 和 $A$ 间能差一个任意常数）。在动作数大、很多状态下动作选择不敏感的任务（Atari 大部分游戏）效果显著。

（这个架构和后面 §Linear FA 的 advantage 概念也呼应：$A^\pi(s,a) = Q^\pi(s,a) - V^\pi(s)$，TD 误差 $\delta$ 是它的无偏估计。）

\paragraph{(6) 熵正则化：在 loss 里鼓励策略保持随机}

\textbf{动机}：deterministic 策略 → 探索停止 → 容易陷入局部最优。

\textbf{做法}：在 loss 里加 $-\beta\, H(\pi(\cdot|s))$，其中 $H(\pi) = -\sum_a \pi(a|s)\log\pi(a|s)$ 是策略熵。$\beta$ 控制随机性强度。

\begin{itemize}
\item \textbf{A2C/PPO}：加一个小系数（$\beta \sim 0.01$）的熵 bonus 防止过早收敛
\item \textbf{SAC（Soft Actor-Critic）}：把它做到极致——目标是最大化 \textbf{奖励 + $\alpha$·策略熵}，称为"最大熵 RL"。这等价于把策略当作 Bellman 软算子的不动点；理论上很优雅
\end{itemize}

熵正则化的副作用：最优策略不是原 MDP 的最优策略（多了一个熵项），但在 $\beta \to 0$ 极限下重合。SAC 还有 auto-tuning $\alpha$ 的版本。

\paragraph{(7) Hindsight Experience Replay（HER）：稀疏奖励的救星}

稀疏奖励（只在最终成功时才有 $+1$，其他全 $0$）是 RL 的克星——绝大多数 trajectory 都拿 $0$ 奖励，几乎学不到东西。

\textbf{HER（Andrychowicz et al.\ 2017）思路}：失败的轨迹也可以当作"达成了另一个目标"的成功轨迹。

\textbf{具体做法}（goal-conditioned 设定下）：
\begin{itemize}
\item 原目标 $g$，trajectory 实际终点 $s_T \neq g$ → 失败，奖励全 $0$
\item Relabel：把目标换成 $g' = s_T$，重新算奖励 → 这个轨迹"达到了"$g'$，变成成功样本！
\item 把 relabeled 轨迹塞进 buffer，正常 off-policy Q-learning
\end{itemize}

\textbf{效果}：在机械臂抓取、push、导航等 goal-conditioned 任务上，从"完全学不到"变成"几小时收敛"。

直觉：每次失败都教会你"如何到达 $s_T$"，哪怕 $s_T$ 不是你想去的地方——总有一天 $s_T$ 可能就是某个用户的目标。把环境反馈的信息密度极大化。

\paragraph{(8) Potential-based reward shaping：加先验知识但不毁最优性}

工程上经常想给奖励函数加点 "提示"——比如"靠近目标更好"。但随便加 shaping 会改变最优策略，得不偿失。

\textbf{Potential-based shaping（Ng, Harada \& Russell 1999）}：选一个势函数 $\Phi: S \to \mathbb{R}$，把奖励改成
\[
R'(s, a, s') = R(s, a) + \gamma \Phi(s') - \Phi(s)
\]
即"势的变化"。

\begin{theorem}[Ng et al.\ 1999]
Potential-based shaping 后的 MDP 与原 MDP 有\textbf{相同的最优策略}。

证明思路：累加 shaping 项的话，中间项全消掉，只剩 $\gamma^T \Phi(s_T) - \Phi(s_0)$，这是个常数（对所有策略一样），所以最优性不变。但每步的"局部信号"更密集了，学习更快。
\end{theorem}

\textbf{用法}：把先验知识"哪些状态更接近成功"编码进 $\Phi$——例如导航任务设 $\Phi(s) = -\|s - s_{\text{goal}}\|$。\textbf{保证最优策略不变，但学习速度可能数量级提升。}

\paragraph{(9) 其他点到为止}

\begin{itemize}
\item \textbf{Distributional RL（C51, QR-DQN, IQN）}：学的不是 $\mathbb E[G_t]$ 而是 $G_t$ 的整个分布。直觉：分布信息比期望多。Atari 上 SOTA 经常用
\item \textbf{NoisyNet}：把网络参数加可学习噪声，作为 $\epsilon$-greedy 的替代——探索量随训练自动衰减
\item \textbf{Frame stacking}：Atari 输入是单帧图像（无速度信息），堆叠最近 4 帧 → 网络能感知运动方向。POMDP 的简单粗暴对策
\item \textbf{Polyak averaging of network weights}：保存训练过程中权重的 EMA 副本作为最终评估模型——通常比最后一步的权重稳定
\item \textbf{Spectral / Layer normalization}：稳定深度网络训练，layer norm 在 RL 中经常比 batch norm 好（batch norm 假设 i.i.d.，replay 出来的 mini-batch 不严格满足）
\end{itemize}

\begin{plainwords}[挑技巧的原则]
不要一上来就把所有技巧都堆上。建议顺序：
\begin{enumerate}
\item 先把 baseline 跑通（Q-learning + replay + target net）
\item 看 loss/Q 是否爆炸 → 加 gradient clipping、Huber
\item 看奖励量级 → 加 normalization
\item 看是否陷入局部最优 → 加熵正则化或更聪明的探索
\item 看是否过估计 → 加 Double
\item 看是否动作不敏感 → 加 Dueling
\item 看是否稀疏奖励 → HER 或 shaping
\item $n$-step、distributional 是锦上添花，先调通基础再上
\end{enumerate}

\textbf{Deep RL 实战很大程度上就是诊断 + 加这些 trick 的循环}。理论保收敛，工程保不发散。
\end{plainwords}

\subsection{算法地图：把这一切收进一张表}

\begin{keypoint}[三维分类总览]
\begin{center}
\small
\begin{tabular}{l|c|c|c|l}
\toprule
算法 & Model & On/off & On/offline & 关键技巧 \\
\midrule
值迭代 & based & — & 不交互 & 直接用$P, R$ \\
策略迭代 & based & — & 不交互 & 评估 + 改进交替 \\
MC评估 & free & on & online & 整段$G_t$取均值 \\
MC control & free & on & online & $\epsilon$-greedy \\
UCB-VI & based & on & online & 估模型 + 乐观bonus \\
SARSA & free & on & online & $\epsilon$-greedy + GLIE \\
Q-learning & free & off & online & $\max$算子 \\
\midrule
DQN & free & off & online & 上面三件套（replay + target net + Double） \\
SAC & free & off & online & + 最大熵 + 双Critic \\
PPO & free & on & online & 信赖域 / clipping \\
CQL / IQL / BCQ & free & off & \textbf{offline} & Q值正则 / 避免max / 行为约束 \\
\bottomrule
\end{tabular}
\end{center}

\textbf{读法}：
\begin{itemize}
\item Model-based通常占offline好处（不用一直交互），但要假设环境可模拟或模型可学
\item On-policy算法（SARSA, PPO）天生不能用replay，只能用最近收集的数据
\item Offline强制off-policy，且需要特殊正则避免OOD问题
\item \textbf{Q-learning家族}（Q-learning $\to$ DQN $\to$ SAC $\to$ CQL）是一条主线——走off-policy路线，逐渐加上replay、target net、Double、最大熵、保守正则等技巧应对不同场景
\end{itemize}

后面§Q-learning收敛性还会从理论角度再补一刀（model-based/free的精确分类、深一些的online/offline讨论），与本节互补——本节关心"工程上是什么"，下一节关心"理论上为什么收敛"。
\end{keypoint}

%=============================================================================
\section{Q-learning收敛性}
%=============================================================================

\begin{plainwords}
值迭代需要知道$P$和$R$。如果不知道呢？

Q-learning：用\textbf{样本}代替期望，在线更新Q值。
\end{plainwords}

\subsection{有模型 vs 无模型：核心区别}

\begin{plainwords}[什么是"模型"？]
在RL中，\textbf{环境模型}指两个东西：
\begin{enumerate}
    \item \textbf{状态转移概率}$P(s'|s,a)$：在状态$s$做动作$a$，跳到$s'$的概率
    \item \textbf{奖励函数}$R(s,a)$：在状态$s$做动作$a$，直接获得多少奖励
\end{enumerate}

\textbf{有模型 vs 无模型}的核心区别：算法更新价值时，是否直接用到$P$和$R$的\textbf{显式数学形式}。
\end{plainwords}

\begin{keypoint}[重要澄清：Model-Free也需要"访问"环境！]
很多人误以为"无模型"意味着"什么都不需要知道"。

\textbf{错！}Model-Free方法（如Q-learning）仍然需要：
\begin{itemize}
    \item 能够在环境中执行动作
    \item 观察环境返回的$(r, s')$
    \item 收集大量交互样本
\end{itemize}

\textbf{它只是不需要知道$P$和$R$的数学公式}，但它需要把环境当作\textbf{黑盒}来查询。

打个比方：
\begin{itemize}
    \item \textbf{Model-Based}：你有一本物理教科书，可以用公式算出球会落在哪
    \item \textbf{Model-Free Online}：你没有教科书，但你可以反复扔球，观察它落在哪
    \item 两者都需要"访问"物理世界，只是方式不同
\end{itemize}
\end{keypoint}

\begin{keypoint}[更精确的三分类]
\begin{center}
\begin{tabular}{l|c|c|c}
\toprule
\textbf{方法} & \textbf{需要$P,R$公式} & \textbf{需要在线交互} & \textbf{例子} \\
\midrule
Model-Based & \textcolor{red}{\ding{51}} & 可选 & 值迭代、MBPO \\
Model-Free Online & \textcolor{green}{\ding{55}} & \textcolor{red}{\ding{51}} & Q-learning、PPO \\
Model-Free Offline & \textcolor{green}{\ding{55}} & \textcolor{green}{\ding{55}} & CQL、IQL、Decision Transformer \\
\bottomrule
\end{tabular}
\end{center}

\textbf{Offline RL}（也叫Batch RL）才是"真正的无模型"：
\begin{itemize}
    \item 不需要$P$和$R$的公式
    \item 不需要和环境交互
    \item 只用预先收集好的历史数据集$\mathcal{D} = \{(s_i, a_i, r_i, s'_i)\}$
\end{itemize}
\end{keypoint}

\begin{example}[三种方法的比喻——学开车]
\textbf{Model-Based}（有驾驶模拟器的物理引擎）：
\begin{itemize}
    \item 你知道方向盘转30度，车会怎么转弯（有$P$的公式）
    \item 可以在脑子里规划："如果我这样打方向盘，车会..."
\end{itemize}

\textbf{Model-Free Online}（真车上路练）：
\begin{itemize}
    \item 你不知道转弯的物理公式
    \item 但你可以反复尝试，每次转完看看效果，慢慢学会
    \item \textbf{需要一辆真车和一条路}
\end{itemize}

\textbf{Model-Free Offline}（看别人的行车记录仪）：
\begin{itemize}
    \item 你不知道物理公式
    \item 你也没有车可以练
    \item 你只有一堆别人开车的视频，从中学习"什么情况该怎么做"
    \item \textbf{这是最难的}：你不能尝试视频里没有的动作
\end{itemize}
\end{example}

\begin{plainwords}[为什么Offline RL很难？]
想象你只能看别人开车的视频学开车，不能自己上手：

\textbf{分布偏移（Distribution Shift）}：
\begin{itemize}
    \item 视频里的司机都是老司机，从不犯错
    \item 你学到的策略只在"正常情况"下有效
    \item 一旦遇到视频里没见过的情况（比如你自己犯了个小错），你就不知道怎么纠正
\end{itemize}

\textbf{过度乐观（Overestimation）}：
\begin{itemize}
    \item 数据集里没有尝试过的动作，Q值可能被高估
    \item 因为你从来没见过那个动作的后果
\end{itemize}

这就是为什么Offline RL需要特殊算法（CQL、IQL等）来处理这些问题。
\end{plainwords}

\begin{keypoint}[值迭代——上帝视角]
值迭代的更新公式：
\[
V_{k+1}(s) = \max_a \sum_{s'} \colorbox{yellow!30}{$P(s'|s,a)$} \left[\colorbox{yellow!30}{$R(s,a)$} + \gamma V_k(s')\right]
\]

\textbf{为什么是有模型的？}
\begin{itemize}
    \item \textbf{显式依赖$P$}：公式里有$\sum_{s'}$和$P(s'|s,a)$，必须预先知道所有可能的下一状态及其概率
    \item \textbf{计算期望}：不是"试一试"，而是直接用概率算出期望值
    \item \textbf{不与环境交互}：只要给它规则书（$P$和$R$），它在脑子里就能算出最优解
\end{itemize}

\textbf{比喻}：就像下棋，你知道规则，可以在脑海里推演，不需要真的落子。
\end{keypoint}

\begin{keypoint}[Q-learning——试错学习（Online Model-Free）]
Q-learning的更新公式：
\[
Q(s,a) \leftarrow Q(s,a) + \alpha\left[\underbrace{r + \gamma \max_{a'} Q(s', a')}_{\text{目标值（来自实际经验）}} - Q(s,a)\right]
\]

\textbf{为什么是无模型的？}
\begin{itemize}
    \item \textbf{$P$和$R$消失了}：公式里没有$\sum_{s'}$，没有$P(s'|s,a)$
    \item \textbf{基于采样}：智能体实际执行动作$a$，环境反馈"你到了$s'$，奖励$r$"
    \item \textbf{依靠交互}：不需要知道"往右走滑倒概率10\%"，只需要摔几次就知道了
\end{itemize}

\textbf{但它仍然需要}：能够和环境交互，收集$(s, a, r, s')$样本！
\end{keypoint}

\begin{example}[计算期望 vs 采样——全班平均身高]
\textbf{值迭代（有模型）}：

老师手里有一张表（模型）：1.6m的学生占10\%，1.7m的占20\%...

直接用公式计算：$1.6 \times 0.1 + 1.7 \times 0.2 + \cdots = \text{期望身高}$

\textbf{不需要看具体学生，只要有概率表就行。}

\textbf{Q-learning（无模型）}：

老师手里没表。随机抓一个学生量身高（采样），1.65m，更新估计值；

再抓一个，1.75m，再更新...

\textbf{不需要概率表，需要一个个具体的样本。}

抓的学生足够多，估计值就逼近真实平均身高。
\end{example}

\begin{center}
\begin{tabular}{l|c|c}
\toprule
\textbf{特性} & \textbf{值迭代} & \textbf{Q-learning} \\
\midrule
类型 & Model-Based（动态规划） & Model-Free（时序差分） \\
核心机制 & 计算期望（Expectation） & 采样（Sampling） \\
是否需要$P(s'|s,a)$ & 是（必须已知） & 否（它是未知黑盒） \\
数据来源 & 环境的数学定义 & 实际交互的经验样本 \\
比喻 & 看着地图规划路线 & 在迷宫里瞎逛摸索路线 \\
\bottomrule
\end{tabular}
\end{center}

\subsection{Q-learning算法}

\begin{example}[Q-learning的生活比喻——试吃学习]
你搬到新城市，想找好吃的餐厅。

\textbf{问题}：你不知道每家店多好吃（不知道$R$），也不知道吃完会想去哪（不知道$P$）。

\textbf{Q-learning策略}：
\begin{enumerate}
    \item 随便选一家店吃（探索）
    \item 吃完打个分（观察奖励$r$）
    \item 更新你对这家店的评分：
    \[
    Q(\text{这家店}) \leftarrow Q(\text{这家店}) + \alpha \times (\text{实际体验} - \text{之前预期})
    \]
\end{enumerate}

\textbf{关键}：你不需要知道"这家店平均多好吃"，只需要\textbf{每次吃完更新一点点}！

吃得多了，你的评分$Q$自然会收敛到真实水平。
\end{example}

\begin{algorithm}[H]
\caption{Q-learning}
\begin{algorithmic}[1]
\State 初始化$Q(s,a)$任意
\For{每步$t$}
    \State 在状态$s_t$选择动作$a_t$（$\epsilon$-greedy）
    \State 观察$r_t$和$s_{t+1}$
    \State $Q(s_t, a_t) \leftarrow Q(s_t, a_t) + \alpha_t \left[r_t + \gamma \max_{a'} Q(s_{t+1}, a') - Q(s_t, a_t)\right]$
\EndFor
\end{algorithmic}
\end{algorithm}

\begin{example}[Q-learning更新的直觉——修正预期]
你以为某餐厅8分（$Q = 8$），结果吃完发现：
\begin{itemize}
    \item 这顿饭7分（$r = 7$）
    \item 吃完心情好，下一个状态价值10分（$\max_{a'} Q(s', a') = 10$）
    \item $\gamma = 0.9$
\end{itemize}

\textbf{TD目标}：$r + \gamma \max Q' = 7 + 0.9 \times 10 = 16$

\textbf{TD误差}：$16 - 8 = 8$（实际比预期好！）

\textbf{更新}（$\alpha = 0.1$）：$Q \leftarrow 8 + 0.1 \times 8 = 8.8$

你把评分从8上调到8.8。多吃几次，评分会稳定在真实水平。
\end{example}

\begin{theorem}[Q-learning收敛 - Watkins \& Dayan, 1992; Jaakkola et al., 1994]
如果：
\begin{enumerate}
    \item 所有$(s,a)$被无限次访问
    \item 学习率满足：$\sum_t \alpha_t(s,a) = \infty$且$\sum_t \alpha_t(s,a)^2 < \infty$
\end{enumerate}
则$Q_t \to Q^*$以概率1。
\end{theorem}

\begin{example}[学习率条件的直觉]
条件1：$\sum_t \alpha_t = \infty$（学习率之和发散）

\textbf{意思}：你要\textbf{持续学习}，不能太早停下来。

如果$\alpha_t$衰减太快（比如$\alpha_t = 1/t^2$），总学习量有限，可能还没学到真值就停了。

条件2：$\sum_t \alpha_t^2 < \infty$（学习率平方之和收敛）

\textbf{意思}：你要\textbf{逐渐冷静}，不能总是大幅调整。

如果$\alpha_t$不衰减（比如恒等于1），噪声会一直影响你，永远不稳定。

\textbf{黄金法则}：$\alpha_t = 1/t$刚好满足两个条件！
\end{example}

\begin{proof}
\textbf{Step 1：写成随机逼近形式}

Q-learning更新可以写成：
\[
Q_{t+1}(s,a) = Q_t(s,a) + \alpha_t(s,a) \cdot \mathbf{1}_{(s_t,a_t)=(s,a)} \cdot [r_t + \gamma \max_{a'} Q_t(s_{t+1}, a') - Q_t(s,a)]
\]

定义：
\begin{itemize}
    \item $F_t(s,a) = R(s,a) + \gamma \sum_{s'} P(s'|s,a) \max_{a'} Q_t(s',a') - Q_t(s,a)$（期望更新方向）
    \item $w_t(s,a) = [r_t + \gamma \max_{a'} Q_t(s_{t+1},a')] - [R(s,a) + \gamma \sum_{s'} P(s'|s,a) \max_{a'} Q_t(s',a')]$（噪声）
\end{itemize}

则更新变成：
\[
Q_{t+1}(s,a) = Q_t(s,a) + \alpha_t \cdot \mathbf{1}_{(s_t,a_t)=(s,a)} \cdot [F_t(s,a) + w_t(s,a)]
\]

\textbf{Step 2：噪声是鞅差}

条件于历史$\mathcal{F}_t = \sigma(s_0, a_0, r_0, \ldots, s_t, a_t)$：
\[
\mathbb{E}[w_t(s,a) | \mathcal{F}_t] = 0
\]
因为$r_t$和$s_{t+1}$的期望正好是$R(s,a)$和$\sum_{s'} P(s'|s,a)(\cdot)$。

\textbf{Step 3：$F_t$具有收缩性质}

令$\Delta_t = Q_t - Q^*$。由于$Q^* = T^* Q^*$（Bellman最优方程）：
\begin{align*}
F_t(s,a) &= (T^* Q_t)(s,a) - Q_t(s,a) \\
&= [(T^* Q_t)(s,a) - (T^* Q^*)(s,a)] + [Q^*(s,a) - Q_t(s,a)] \\
&= [(T^* Q_t)(s,a) - (T^* Q^*)(s,a)] - \Delta_t(s,a)
\end{align*}

由$T^*$的收缩性（前面证明过）：
\[
|(T^* Q_t)(s,a) - (T^* Q^*)(s,a)| \leq \gamma \|\Delta_t\|_\infty
\]

所以：
\[
|F_t(s,a) + \Delta_t(s,a)| \leq \gamma \|\Delta_t\|_\infty
\]

\textbf{Step 4：应用随机逼近定理}

我们需要Jaakkola等人的随机逼近定理：

\begin{theorem}[随机逼近收敛条件]
考虑迭代：$x_{t+1} = x_t + \alpha_t(F_t(x_t) + w_t)$

如果：
\begin{enumerate}
    \item $\mathbb{E}[w_t | \mathcal{F}_t] = 0$，$\mathbb{E}[w_t^2 | \mathcal{F}_t] \leq C(1 + \|x_t\|^2)$
    \item 存在$\gamma < 1$使得$\|F_t(x)\| \leq \gamma \|x\|$（收缩性）
    \item $\sum_t \alpha_t = \infty$，$\sum_t \alpha_t^2 < \infty$
\end{enumerate}
则$x_t \to 0$ w.p.1。
\end{theorem}

对Q-learning，令$x_t = \Delta_t = Q_t - Q^*$：

\begin{itemize}
    \item 条件1：噪声$w_t$有界方差（因为$Q_t$有界，$r_t$有界）\checkmark
    \item 条件2：$F_t(\Delta) = (T^* Q) - Q = (T^*(Q^* + \Delta)) - (Q^* + \Delta)$满足$|F_t + \Delta| \leq \gamma\|\Delta\|$ \checkmark
    \item 条件3：学习率条件 \checkmark
\end{itemize}

因此$\Delta_t \to 0$，即$Q_t \to Q^*$。

\textbf{Step 5：处理异步更新}

上面假设每步更新所有$(s,a)$。实际Q-learning只更新访问到的$(s,a)$。

关键：如果每个$(s,a)$被无限次访问，它的有效学习率序列$\{\alpha_{t_i}\}$仍满足条件3。

由独立的随机逼近定理，每个$(s,a)$的$Q$值都收敛到$Q^*$。
\end{proof}

\begin{calculation}[Q-learning收敛速度]
理论：$\mathbb{E}[\|\Delta_t\|^2] = O(1/t)$（使用$\alpha_t = 1/t$）

实际例子：$|S|=100$，$|A|=4$，$\gamma=0.9$

要达到$\|\Delta\| \leq 0.1$：
\begin{itemize}
    \item 需要每个$(s,a)$被更新约$O(1/0.1^2) = 100$次
    \item 总共$100 \times 4 \times 100 = 40000$次更新
    \item 如果探索均匀，需要约40000次转移
\end{itemize}

实际中由于探索不均匀，可能需要更多。
\end{calculation}

\begin{keypoint}[Q-learning vs 值迭代]
\begin{center}
\begin{tabular}{l|c|c}
\toprule
& 值迭代 & Q-learning \\
\midrule
需要模型 & 是 & 否 \\
更新方式 & 同步（所有状态） & 异步（访问到的） \\
收敛速度 & $O(\gamma^k)$（线性） & $O(1/k)$（次线性） \\
每次代价 & $O(|S|^2|A|)$ & $O(1)$ \\
\bottomrule
\end{tabular}
\end{center}

Q-learning牺牲了收敛速度，换来了不需要模型的优势。
\end{keypoint}

%=============================================================================
\section{线性函数逼近的收敛性}
%=============================================================================

\begin{plainwords}
当$|S|$很大（比如$10^6$），不能存储完整的$V(s)$表。

解决方案：\textbf{函数逼近}$V_\theta(s) = \theta^T \phi(s)$

问题：还能收敛吗？
\end{plainwords}

\subsection{Linear TD(0)}

\begin{definition}[Linear TD(0)]
\[
\theta_{t+1} = \theta_t + \alpha_t \underbrace{(r_t + \gamma \theta_t^T \phi(s_{t+1}) - \theta_t^T \phi(s_t))}_{\text{TD误差 } \delta_t} \phi(s_t)
\]
\end{definition}

\begin{theorem}[Linear TD(0)收敛 - Tsitsiklis \& Van Roy, 1997]
在on-policy采样下，Linear TD(0)收敛到\textbf{投影Bellman方程}的解：
\[
\Phi \theta^* = \Pi_D T^\pi (\Phi \theta^*)
\]
其中$\Pi_D$是在状态分布$D$加权下到特征空间的投影。
\end{theorem}

\begin{plainwords}[投影Bellman方程]
真实$V^\pi$可能不在线性空间里。TD(0)找的是：在\textbf{能表示的函数中}，最接近满足Bellman方程的那个。
\end{plainwords}

\begin{proof}
\textbf{Step 1：期望更新}

$\theta$的期望更新方向：
\[
\mathbb{E}[\delta_t \phi(s_t)] = \Phi^T D (T^\pi \Phi \theta - \Phi \theta)
\]

\textbf{Step 2：不动点}

令期望更新为零：$\Phi^T D (T^\pi \Phi \theta^* - \Phi \theta^*) = 0$

这意味着$T^\pi \Phi\theta^* - \Phi\theta^* \perp \text{col}(\Phi)$在$D$-内积下，即$\Phi\theta^* = \Pi_D T^\pi \Phi\theta^*$。

\textbf{Step 3：$\Pi_D T^\pi$是收缩}

$\|\Pi_D T^\pi V - \Pi_D T^\pi V'\|_D \leq \|T^\pi V - T^\pi V'\|_D \leq \gamma\|V-V'\|_D$

由Banach定理，$\theta^*$存在且唯一。
\end{proof}

\begin{theorem}[逼近误差界]
\[
\|V_{\theta^*} - V^\pi\|_D \leq \frac{1}{\sqrt{1-\gamma^2}} \min_\theta \|V^\pi - \Phi\theta\|_D
\]
\end{theorem}

\begin{calculation}[误差放大]
\begin{center}
\begin{tabular}{c|c|l}
\toprule
$\gamma$ & 放大因子$\frac{1}{\sqrt{1-\gamma^2}}$ & 含义 \\
\midrule
0.9 & 2.3 & 最佳逼近误差放大2.3倍 \\
0.99 & 7.1 & 放大7倍 \\
0.999 & 22.4 & 放大22倍 \\
\bottomrule
\end{tabular}
\end{center}

$\gamma$越大，线性逼近误差放大越严重！
\end{calculation}

\subsection{Deadly Triad：发散的危险}

\begin{keypoint}[致命三角]
以下三者组合可能导致\textbf{发散}：
\begin{enumerate}
    \item \textbf{函数逼近}（不是表格）
    \item \textbf{Bootstrapping}（用估计更新估计）
    \item \textbf{Off-policy}（行为策略$\neq$目标策略）
\end{enumerate}
\end{keypoint}

\begin{example}[Deadly Triad的生活比喻——谣言传播]
想象一个办公室里传谣言：

\textbf{函数逼近}：你不认识所有人，只能根据"职位相似的人想法差不多"来猜测。

\textbf{Bootstrapping}：你从别人那听到消息，就信以为真，然后传给下一个人。

\textbf{Off-policy}：你听的是A部门的消息，但你在B部门传。

\textbf{单独一个问题不大}：
\begin{itemize}
    \item 只有函数逼近：你直接问每个人，虽然有猜测但能纠正
    \item 只有bootstrapping：消息来源准确，传着传着会稳定
    \item 只有off-policy：至少你传的是真实消息
\end{itemize}

\textbf{三个一起就炸了}：
\begin{itemize}
    \item 你猜B部门的人和A部门某人想法一样
    \item 你从A部门听到一个猜测（不是事实）
    \item 你把这个猜测传给B部门，B部门的人又传回A部门
    \item 误差越传越大，最后谣言离谱到天上去！
\end{itemize}
\end{example}

\begin{example}[Baird反例]
Baird (1995)构造了一个7状态MDP，使用Linear Q-learning（off-policy），参数$\theta \to \infty$！

原因：off-policy采样分布与更新目标不匹配，更新方向不是下降方向。
\end{example}

\begin{theorem}[Gradient TD收敛 - Sutton et al., 2009]
GTD2算法通过\textbf{显式最小化}投影Bellman误差，在off-policy下也收敛。
\end{theorem}

\begin{plainwords}[如何避免Deadly Triad]
\begin{enumerate}
    \item \textbf{用表格}：如果状态空间小，直接用表格，不用函数逼近
    \item \textbf{用Monte Carlo}：不用bootstrapping，等episode结束再更新
    \item \textbf{用on-policy}：只学习你实际执行的策略（如SARSA）
    \item \textbf{用GTD}：如果必须三者都要，用Gradient TD类方法
\end{enumerate}
\end{plainwords}

%=============================================================================
\section{两个关键引理}
%=============================================================================

\begin{plainwords}
接下来介绍两个在RL理论中反复出现的引理，它们是后续很多结果的基础。
\end{plainwords}

\subsection{Simulation Lemma：模型误差如何传播}

\begin{plainwords}
Model-based RL学习环境模型$\hat{P}$，然后在模型上规划。

问题：模型不完美，误差如何影响最终策略？
\end{plainwords}

\begin{example}[Simulation Lemma的生活版——用地图导航]
你要从A到B，用的是一张旧地图（模型$\hat{P}$）。

\textbf{问题}：地图和实际道路（真实$P$）有些出入，会导致多大问题？

\textbf{直觉分析}：
\begin{itemize}
    \item 如果只走1步，地图误差$\epsilon$就是实际误差$\epsilon$
    \item 如果走2步，第1步的误差让你到错误位置，第2步在错误位置又犯错
    \item 走$n$步，误差\textbf{累积}！
\end{itemize}

\textbf{更糟糕的是}：未来的误差要\textbf{折现回今天}。

如果$\gamma = 0.99$（很看重未来），折现因子是$1/(1-0.99)^2 = 10000$！

\textbf{结论}：地图误差1\%，导航效果可能差100倍！
\end{example}

\begin{theorem}[Simulation Lemma]
设真实MDP为$(P, R)$，学习到的模型为$(\hat{P}, \hat{R})$，$\hat{\pi}$是模型上的最优策略。

则在真实环境中：
\[
V^*(s) - V^{\hat{\pi}}(s) \leq \frac{2}{(1-\gamma)^2} \left(\epsilon_R + \gamma \epsilon_P \cdot \frac{R_{max}}{1-\gamma}\right)
\]
其中$\epsilon_R = \max_{s,a}|R - \hat{R}|$，$\epsilon_P = \max_{s,a}\|P - \hat{P}\|_1$。
\end{theorem}

\begin{example}[误差放大的数值感受]
$\gamma = 0.99$，$R_{max} = 1$，模型误差$\epsilon_P = \epsilon_R = 0.01$（1\%）

放大因子：$\frac{2}{(1-0.99)^2} = \frac{2}{0.0001} = 20000$

策略次优性上界：$20000 \times (0.01 + 0.99 \times 0.01 \times 100) \approx 20000$

\textbf{1\%的模型误差，策略可能差20000分}！

这就是为什么model-based RL对模型精度要求极高。
\end{example}

\begin{proof}
\textbf{Step 1}：分解误差
\[
V^* - V^{\hat{\pi}} = \underbrace{(V^* - \hat{V}^*)}_{\text{模型乐观/悲观}} + \underbrace{(\hat{V}^{\hat{\pi}} - V^{\hat{\pi}})}_{\text{策略在模型vs真实}}
\]

\textbf{Step 2}：界定单策略误差

对任意策略$\pi$，比较$V^\pi$（真实）和$\hat{V}^\pi$（模型）：
\[
|V^\pi(s) - \hat{V}^\pi(s)| \leq \frac{\epsilon_R}{1-\gamma} + \frac{\gamma \epsilon_P R_{max}}{(1-\gamma)^2}
\]

\textbf{Step 3}：合并两项。
\end{proof}

\subsection{为什么Model-Based RL在实践中常常不如Model-Free？}

\begin{plainwords}[一个尖锐的观察]
理论上，Model-Based RL应该更好：样本效率高、能规划、懂因果。

但实践中，Model-Free方法（PPO, SAC, DQN）占据绝对主导：ChatGPT的RLHF、OpenAI Five打Dota、机器人控制...

为什么？因为MBRL有四大痛点，让它\textbf{"一顿操作猛如虎，一看效果二百五"}。
\end{plainwords}

\begin{keypoint}[痛点1：复合误差——最致命的杀手]
这就是Simulation Lemma警告我们的！

\textbf{场景}：你的模型精度99\%（已经很高了），需要规划100步。

\begin{center}
\begin{tabular}{c|c}
\toprule
步数 & 累积精度 \\
\midrule
1 & $0.99$ \\
10 & $0.99^{10} = 0.90$ \\
50 & $0.99^{50} = 0.61$ \\
100 & $0.99^{100} = 0.36$ \\
\bottomrule
\end{tabular}
\end{center}

\textbf{结果}：智能体在脑海里规划出完美路径，现实中走到第10步就掉沟里了。

\textbf{Model-Free没这问题}：它不预测未来，只记住"这个状态做这个动作大概率是好的"。不需要多步推演，没有误差累积。
\end{keypoint}

\begin{keypoint}[痛点2：Vapnik原则——不要解决更难的问题]
统计学习大佬Vapnik说：\textbf{"尽量直接解决问题，而不要通过解决一个更难的问题来解决它。"}

\textbf{MBRL的做法}（更难的问题）：
\begin{itemize}
    \item 试图学会整个世界的物理规律$P(s'|s,a)$
    \item 为了学会打网球，先学空气动力学、重力、肌肉力学，然后在脑子里模拟
    \item 这太难了！
\end{itemize}

\textbf{MFRL的做法}（直接的问题）：
\begin{itemize}
    \item 试图学会怎么赢（策略$\pi(a|s)$）
    \item 不管空气动力学，只知道"球飞过来，挥拍就能赢"
\end{itemize}

\textbf{现实}：策略往往很简单（"见到悬崖就左转"），但环境动力学极其复杂（悬崖边的风速、地面摩擦力）。非要建模那个复杂环境，属于\textbf{舍近求远}。
\end{keypoint}

\begin{keypoint}[痛点3：目标不匹配——预测准$\neq$拿分高]
训练环境模型用的是MSE损失：让预测画面和真实画面像素越接近越好。

\textbf{例子}：玩《超级马里奥》

模型的重点：屏幕背景有大量云朵和树叶，模型花90\%力气预测云朵怎么飘（像素多，MSE权重高）

被忽略的细节：屏幕角落有一个只有几像素的\textbf{子弹}飞过来。模型觉得预测这个子弹误差才几个像素，无所谓。

\textbf{后果}：模型能完美生成背景风景画，但完全忽略了致命子弹。智能体觉得前面安全，一走过去就挂了。

这就是：\textbf{Model预测得准（Loss低），但对RL任务（Reward）没帮助}。
\end{keypoint}

\begin{keypoint}[痛点4：幻觉作弊——智能体很鸡贼]
智能体会\textbf{利用模型的Bug刷分}，而不是学会真正的技能。

\textbf{案例}：模拟器里墙壁碰撞检测有点Bug。

智能体发现：以某个奇怪角度撞墙，能瞬移或获得无限分！

\textbf{结果}：
\begin{itemize}
    \item 在模型里：表现得像神，拿到正无穷奖励
    \item 部署到真实环境：只会傻傻往墙上撞，极其弱智
\end{itemize}

这就是为什么Sim-to-Real（仿真到现实）如此困难。
\end{keypoint}

\begin{example}[MBRL的复活之路——隐空间模型]
既然预测像素（真实物理世界）太难，那就\textbf{不预测像素}！

\textbf{MuZero / AlphaZero}：
\begin{itemize}
    \item 不预测棋盘画面，而是预测"价值"和"策略"
    \item 学会了一个"抽象的规则"
    \item 在棋类和Atari游戏中封神
\end{itemize}

\textbf{Dreamer (V1/V2/V3)}：
\begin{itemize}
    \item 不预测具体像素，而是在压缩的Latent Space（潜空间）里"做梦"
    \item 大大减少了复合误差
\end{itemize}

\textbf{结论}：
\begin{itemize}
    \item "显式"建模物理世界（Old School MBRL）→ 不太行，世界太复杂
    \item "隐式"建模（Latent MBRL）→ 正在成为SOTA，但需要极高算法技巧
    \item 目前PPO这种"傻瓜式"Model-Free算法仍然最好用、最稳定
\end{itemize}
\end{example}

\begin{center}
\begin{tabular}{l|c|c}
\toprule
\textbf{对比} & \textbf{Model-Based} & \textbf{Model-Free} \\
\midrule
样本效率 & 高（理论上） & 低 \\
复合误差 & 致命问题 & 无 \\
实现难度 & 高 & 低 \\
稳定性 & 差 & 好 \\
工业应用 & 少 & 主流（PPO, SAC） \\
\bottomrule
\end{tabular}
\end{center}

\subsection{Performance Difference Lemma}

\begin{theorem}[Performance Difference Lemma - Kakade \& Langford, 2002]
\[
V^{\pi'}(s) - V^\pi(s) = \frac{1}{1-\gamma} \mathbb{E}_{s' \sim d^{\pi'}_s}\left[\mathbb{E}_{a \sim \pi'}[A^\pi(s', a)]\right]
\]
其中$A^\pi(s,a) = Q^\pi(s,a) - V^\pi(s)$是优势函数。
\end{theorem}

\begin{example}[Performance Difference的生活版——换工作]
你现在的工作策略是$\pi$（比如"遇到问题就加班"）。

你考虑新策略$\pi'$（比如"遇到问题先求助同事"）。

\textbf{问}：新策略比旧策略好多少？

\textbf{优势函数}$A^\pi(s, a)$的意思：
\begin{itemize}
    \item 在状态$s$（比如"遇到bug"）
    \item 如果用旧策略，期望收益是$V^\pi(s)$（加班解决，累但有效）
    \item 如果这次选动作$a$（求助同事），期望收益是$Q^\pi(s, a)$
    \item $A^\pi(s, a) = Q^\pi(s, a) - V^\pi(s)$就是"这次选$a$比按旧习惯\textbf{好多少}"
\end{itemize}

\textbf{Performance Difference Lemma说}：

新策略的总收益 - 旧策略的总收益 = 新策略访问的所有状态上，新策略的动作比旧策略\textbf{平均好多少}。

\textbf{关键洞察}：用\textbf{旧策略的}$A^\pi$评估\textbf{新策略}$\pi'$！
\end{example}

\begin{proof}
\textbf{Step 1}：展开$V^{\pi'}$并加减$V^\pi$
\begin{align*}
V^{\pi'}(s) &= \mathbb{E}_{\pi'}\left[\sum_{t=0}^\infty \gamma^t r_t\right] \\
&= \mathbb{E}_{\pi'}\left[\sum_{t=0}^\infty \gamma^t (r_t + V^\pi(s_t) - V^\pi(s_t))\right] \\
&= \mathbb{E}_{\pi'}\left[\sum_{t=0}^\infty \gamma^t (r_t + \gamma V^\pi(s_{t+1}) - V^\pi(s_t))\right] + V^\pi(s)
\end{align*}
最后用了telescoping：$\sum_t \gamma^t(V^\pi(s_t) - \gamma V^\pi(s_{t+1})) = V^\pi(s_0)$。

\textbf{Step 2}：识别优势函数

$r_t + \gamma V^\pi(s_{t+1}) - V^\pi(s_t) = Q^\pi(s_t, a_t) - V^\pi(s_t) = A^\pi(s_t, a_t)$

所以$V^{\pi'}(s) - V^\pi(s) = \mathbb{E}_{\pi'}[\sum_t \gamma^t A^\pi(s_t, a_t)]$。

\textbf{Step 3}：重写为状态分布的期望。
\end{proof}

\begin{corollary}[策略改进条件]
如果对所有$s$：$\mathbb{E}_{a \sim \pi'}[A^\pi(s, a)] \geq 0$，则$V^{\pi'} \geq V^\pi$。
\end{corollary}

\begin{example}[策略改进的直觉]
如果新策略$\pi'$在\textbf{每个状态}的动作都\textbf{不比}旧策略差（$A^\pi \geq 0$），那新策略整体更好。

这就像：如果你每个决策都至少和以前一样好，那你的人生不会变差！
\end{example}

\begin{keypoint}[Performance Difference Lemma的应用]
\begin{itemize}
    \item \textbf{策略迭代}：贪婪策略保证$A^\pi(s, \pi'(s)) \geq 0$
    \item \textbf{TRPO/PPO}（下两周）：用$A^\pi$构造代理目标，约束$\pi'$不离$\pi$太远
\end{itemize}
\end{keypoint}

%=============================================================================
\section{Multi-Agent RL的收敛性}
%=============================================================================

\begin{plainwords}
当有\textbf{多个agent}同时学习时：
\begin{itemize}
    \item 每个agent的最优策略依赖于其他agent
    \item 其他agent也在学习，环境\textbf{非平稳}
    \item 目标是\textbf{纳什均衡}，不是单方最优
\end{itemize}
\end{plainwords}

\begin{example}[Multi-Agent的生活版——室友博弈]
你和室友共用一个厨房，每天晚上都要决定：

\textbf{动作}：\{早点做饭(6pm), 晚点做饭(8pm)\}

\textbf{收益}：
\begin{itemize}
    \item 两人都6pm：厨房拥挤，各得3分
    \item 两人都8pm：厨房拥挤，各得3分
    \item 一人6pm一人8pm：都舒服，各得5分
\end{itemize}

\textbf{问题}：你应该选什么时间？

\textbf{困难}：你的最优选择\textbf{取决于室友选什么}！
\begin{itemize}
    \item 如果室友选6pm，你应该选8pm
    \item 如果室友选8pm，你应该选6pm
    \item 但室友也在想同样的问题...
\end{itemize}

\textbf{纳什均衡}：双方都随机选（各50\%），期望收益4分。谁都不想单方面改变。
\end{example}

\begin{definition}[随机博弈]
$n$人随机博弈$(S, \{A_i\}, P, \{R_i\}, \gamma)$：每个玩家$i$有自己的动作空间$A_i$和奖励函数$R_i$。
\end{definition}

\begin{definition}[纳什均衡]
$(\pi_1^*, \ldots, \pi_n^*)$是纳什均衡，如果没有玩家能通过单方面改变策略获益。
\end{definition}

\begin{example}[纳什均衡的直觉——交通灯]
两辆车在十字路口相遇，都想先过。

\textbf{策略}：\{冲过去, 让一下\}

\textbf{收益矩阵}（你的收益）：
\begin{center}
\begin{tabular}{c|cc}
& 对方冲 & 对方让 \\
\hline
你冲 & -100（撞车） & 10（你先过） \\
你让 & 5（对方先过） & 0（都让，尴尬）
\end{tabular}
\end{center}

\textbf{纳什均衡}：一个冲一个让。如果你知道对方会让，你就冲；如果你知道对方会冲，你就让。

在这个均衡点，谁都不想单方面改变！
\end{example}

\subsection{零和博弈：Minimax收敛}

\begin{definition}[零和随机博弈]
两人零和随机博弈：$R_1(s, a_1, a_2) = -R_2(s, a_1, a_2)$。

玩家1最大化，玩家2最小化。
\end{definition}

\begin{theorem}[Minimax定理 - von Neumann]
对任意矩阵博弈$M$：
\[
\max_{\pi_1 \in \Delta(A_1)} \min_{a_2 \in A_2} \sum_{a_1} \pi_1(a_1) M(a_1, a_2) = \min_{\pi_2 \in \Delta(A_2)} \max_{a_1 \in A_1} \sum_{a_2} \pi_2(a_2) M(a_1, a_2)
\]
这个共同值记为$\text{val}(M)$。
\end{theorem}

\begin{theorem}[Minimax-Q收敛 - Littman, 1994]
Minimax-Q学习更新：
\[
Q(s, a_1, a_2) \leftarrow (1-\alpha)Q(s,a_1,a_2) + \alpha[r + \gamma \text{val}(Q(s', \cdot, \cdot))]
\]
在满足标准随机逼近条件下，$Q \to Q^*$（纳什均衡Q函数）。
\end{theorem}

\begin{proof}
\textbf{Step 1：定义Minimax Bellman算子}

\[
(T Q)(s, a_1, a_2) = R(s, a_1, a_2) + \gamma \sum_{s'} P(s'|s, a_1, a_2) \text{val}(Q(s', \cdot, \cdot))
\]

\textbf{Step 2：证明$T$是$\gamma$-收缩}

关键引理：$\text{val}(\cdot)$是非扩张的（non-expansive）。

\begin{lemma}
对任意矩阵$M, M'$：$|\text{val}(M) - \text{val}(M')| \leq \|M - M'\|_\infty$
\end{lemma}

\begin{proof}[引理证明]
设$\pi_1^*$是$M$的最优混合策略，$\pi_2^*$是$M'$的最优混合策略。

\begin{align*}
\text{val}(M) &= \max_{\pi_1} \min_{a_2} \pi_1^T M e_{a_2} = \min_{a_2} (\pi_1^*)^T M e_{a_2} \\
&\leq \min_{a_2} (\pi_1^*)^T M' e_{a_2} + \|M - M'\|_\infty \\
&\leq \max_{\pi_1} \min_{a_2} \pi_1^T M' e_{a_2} + \|M - M'\|_\infty \\
&= \text{val}(M') + \|M - M'\|_\infty
\end{align*}

对称地，$\text{val}(M') \leq \text{val}(M) + \|M - M'\|_\infty$。

所以$|\text{val}(M) - \text{val}(M')| \leq \|M - M'\|_\infty$。
\end{proof}

现在证明$T$是收缩：

对任意$Q, Q'$和$(s, a_1, a_2)$：
\begin{align*}
|(TQ)(s,a_1,a_2) - (TQ')(s,a_1,a_2)| &= \gamma \left|\sum_{s'} P(s'|s,a_1,a_2)[\text{val}(Q(s',\cdot,\cdot)) - \text{val}(Q'(s',\cdot,\cdot))]\right| \\
&\leq \gamma \sum_{s'} P(s'|s,a_1,a_2) |\text{val}(Q(s',\cdot,\cdot)) - \text{val}(Q'(s',\cdot,\cdot))| \\
&\leq \gamma \sum_{s'} P(s'|s,a_1,a_2) \|Q(s',\cdot,\cdot) - Q'(s',\cdot,\cdot)\|_\infty \\
&\leq \gamma \|Q - Q'\|_\infty
\end{align*}

\textbf{Step 3：Banach不动点定理}

$T$是$\gamma$-收缩，所以存在唯一不动点$Q^*$，且值迭代收敛。

\textbf{Step 4：随机逼近收敛}

Minimax-Q是$T$的随机逼近版本。设$\Delta_t = Q_t - Q^*$：
\[
\Delta_{t+1}(s,a_1,a_2) = (1-\alpha_t)\Delta_t(s,a_1,a_2) + \alpha_t[\gamma(\text{val}(Q_t) - \text{val}(Q^*)) + w_t]
\]

由$\text{val}$的非扩张性：$|\text{val}(Q_t) - \text{val}(Q^*)| \leq \|\Delta_t\|_\infty$

这是带噪声的收缩迭代，由随机逼近理论收敛。
\end{proof}

\begin{calculation}[石头剪刀布的值]
收益矩阵（玩家1）：
\[
M = \begin{pmatrix} 0 & -1 & 1 \\ 1 & 0 & -1 \\ -1 & 1 & 0 \end{pmatrix}
\]

求$\text{val}(M)$：设玩家1策略$\pi = (p, q, 1-p-q)$。

对手最优响应：选使$\pi^T M e_j$最小的$j$。

$\pi^T M e_1 = q - (1-p-q) = p + 2q - 1$

$\pi^T M e_2 = -p + (1-p-q) = 1 - 2p - q$

$\pi^T M e_3 = p - q$

要最大化$\min_j \pi^T M e_j$，令三者相等：

$p + 2q - 1 = 1 - 2p - q = p - q$

从第一和第三个：$p + 2q - 1 = p - q$，得$3q = 1$，$q = 1/3$。

从第二和第三个：$1 - 2p - q = p - q$，得$1 = 3p$，$p = 1/3$。

所以$\pi^* = (1/3, 1/3, 1/3)$，$\text{val}(M) = 1/3 - 1/3 = 0$。

\textbf{纳什均衡}：双方都用$(1/3, 1/3, 1/3)$，值为0（公平博弈）。
\end{calculation}

\subsection{一般和博弈：收敛困难}

\begin{theorem}[独立Q-learning可能不收敛]
存在一般和博弈，独立Q-learning在策略空间中循环。
\end{theorem}

\begin{example}[Shapley博弈]
\[
\begin{array}{c|ccc}
& A & B & C \\
\hline
A & (0,0) & (1,0) & (0,1) \\
B & (0,1) & (0,0) & (1,0) \\
C & (1,0) & (0,1) & (0,0)
\end{array}
\]

独立Q-learning在$(A,A) \to (B,B) \to (C,C) \to (A,A)$循环，永远不收敛到混合均衡$(1/3, 1/3, 1/3)$。
\end{example}

%=============================================================================
\section{TD($\lambda$)与资格迹}
%=============================================================================

\begin{plainwords}
TD(0)只看\textbf{一步}：$V(s) \leftarrow V(s) + \alpha(r + \gamma V(s') - V(s))$

Monte Carlo看\textbf{整条轨迹}：$V(s) \leftarrow V(s) + \alpha(G_t - V(s))$

TD($\lambda$)用$\lambda \in [0,1]$在两者之间平滑过渡。
\end{plainwords}

\begin{example}[TD(0) vs MC vs TD($\lambda$)——考试成绩预测]
你想预测学生的期末成绩，观察到一个学生的学习轨迹：

第1周作业：80分 → 第2周作业：85分 → 第3周测验：90分 → 期末：88分

\textbf{TD(0)}（只看下一步）：
\begin{itemize}
    \item 预测第1周价值时，只看第2周：$V(\text{第1周}) \approx 80 + \gamma \times V(\text{第2周})$
    \item 问题：如果第2周的预测不准，误差会传播
\end{itemize}

\textbf{Monte Carlo}（看整条轨迹）：
\begin{itemize}
    \item 等到期末出分，用真实累积：$V(\text{第1周}) = 80 + 0.9 \times 85 + 0.81 \times 90 + ...$
    \item 问题：必须等到episode结束，而且方差大
\end{itemize}

\textbf{TD($\lambda$)}（综合考虑）：
\begin{itemize}
    \item 给不同远近的信号加权：近的权重高（更可靠），远的权重低
    \item $\lambda = 0$：退化为TD(0)
    \item $\lambda = 1$：退化为MC
    \item $\lambda = 0.9$：通常效果最好
\end{itemize}
\end{example}

\begin{definition}[$\lambda$-回报]
\[
G_t^\lambda = (1-\lambda) \sum_{n=1}^\infty \lambda^{n-1} G_t^{(n)}
\]
是所有$n$-步回报$G_t^{(n)} = \sum_{i=0}^{n-1} \gamma^i r_{t+i} + \gamma^n V(s_{t+n})$的指数加权平均。
\end{definition}

\begin{example}[$\lambda$-回报的直觉——综合多个消息源]
你想预测明天的天气，有几个信息源：
\begin{itemize}
    \item 今晚的天气预报（1步预测）：80\%准确
    \item 看云彩自己判断（2步预测）：70\%准确
    \item 农历节气（长期预测）：60\%准确
\end{itemize}

\textbf{TD(0)}：只信天气预报

\textbf{MC}：只信农历节气（等到结果出来）

\textbf{TD($\lambda$)}：综合所有信息，近期信息权重高：
\[
\text{预测} = 0.3 \times \text{天气预报} + 0.21 \times \text{云彩} + 0.15 \times \text{节气} + ...
\]
（$\lambda = 0.7$时的权重分配）
\end{example}

\begin{theorem}[前向=后向]
\textbf{前向视角}：更新目标是$G_t^\lambda$

\textbf{后向视角}：用资格迹$e(s) \leftarrow \gamma\lambda e(s) + \mathbf{1}_{s_t=s}$

两者\textbf{等价}！关键是$G_t^\lambda - V(s_t) = \sum_{k=0}^\infty (\gamma\lambda)^k \delta_{t+k}$。
\end{theorem}

\begin{example}[资格迹的直觉——功劳分配]
你的足球队赢了比赛，怎么分配功劳？

\textbf{TD(0)}：只奖励进球的人

\textbf{MC}：平均分配给所有人

\textbf{TD($\lambda$)}（资格迹）：
\begin{itemize}
    \item 刚传球的人，功劳大（资格迹高）
    \item 10分钟前传球的人，功劳小一点（资格迹衰减）
    \item 上半场的助攻，功劳很小（资格迹已经很低）
\end{itemize}

资格迹就是\textbf{记忆}：记住谁最近访问过，用$\gamma\lambda$的速度遗忘。
\end{example}

\begin{keypoint}[$\lambda$的选择]
\begin{center}
\begin{tabular}{c|c|c|l}
\toprule
$\lambda$ & 偏差 & 方差 & 适用场景 \\
\midrule
0（TD(0)） & 高 & 低 & 模型准确，想快速学习 \\
1（MC） & 低 & 高 & Episode短，奖励稀疏 \\
0.9 & 中等 & 中等 & \textbf{通常最佳选择} \\
\bottomrule
\end{tabular}
\end{center}
\end{keypoint}

%=============================================================================
\section{LSTD/LSPI：批量方法}
%=============================================================================

\begin{plainwords}
TD是增量更新。能不能收集一堆数据，一次性算出最优参数？

LSTD：闭式解，无需调学习率！
\end{plainwords}

\begin{example}[LSTD vs TD的比喻——学习骑自行车]
\textbf{TD方法}（在线学习）：
\begin{itemize}
    \item 骑一小段，摔了，调整一点姿势
    \item 再骑一小段，又摔了，再调整一点
    \item 每次只根据\textbf{当前这一跤}调整
    \item 问题：调整幅度（学习率）很难选，太大会乱晃，太小学得慢
\end{itemize}

\textbf{LSTD方法}（批量学习）：
\begin{itemize}
    \item 先随便骑，录下视频，摔100跤
    \item 回家看完所有视频，\textbf{一次性}总结出最佳姿势
    \item 直接用公式算，不需要调学习率
    \item 优点：数据利用效率高；缺点：要等收集完数据
\end{itemize}
\end{example}

\begin{theorem}[LSTD解]
$\theta^* = A^{-1}b$，其中：
\[
A = \mathbb{E}[\phi(s)(\phi(s) - \gamma\phi(s'))^T], \quad b = \mathbb{E}[\phi(s)r]
\]
\end{theorem}

\begin{example}[LSTD公式的直觉——最小二乘]
LSTD本质上是在解一个\textbf{最小二乘问题}：

"找一个线性函数$V_\theta(s) = \theta^T \phi(s)$，使得它尽可能满足Bellman方程"

就像高中学的线性回归：给定数据点，找一条直线最好地拟合它们。

只不过这里的"数据点"是$(s, r + \gamma V(s'))$，我们要拟合的是Bellman方程。
\end{example}

\begin{keypoint}[LSTD的优势]
\begin{enumerate}
    \item \textbf{数据高效}：固定数据集，不需要在线交互
    \item \textbf{无学习率}：直接求解，不用调参
    \item \textbf{Off-policy}：数据可来自任何策略
\end{enumerate}

代价：$O(d^2)$内存，$O(d^3)$计算（$d$是特征维度）。
\end{keypoint}

\begin{example}[什么时候用LSTD]
\textbf{适合LSTD的场景}：
\begin{itemize}
    \item 你有一批历史数据（比如过去的用户行为日志）
    \item 不能在线交互（比如医疗决策，不能随便试）
    \item 特征维度不太高（$d < 1000$）
\end{itemize}

\textbf{不适合LSTD的场景}：
\begin{itemize}
    \item 可以持续收集新数据
    \item 特征维度很高（$d > 10000$，矩阵求逆太慢）
    \item 需要实时更新
\end{itemize}
\end{example}

%=============================================================================
\section{平均奖励MDP}
%=============================================================================

\begin{plainwords}
之前是折扣奖励$\sum_t \gamma^t r_t$。另一种设定：平均奖励$\lim_{T\to\infty} \frac{1}{T}\sum_t r_t$。

对"无限期运行"的系统更自然：服务器调度、库存管理等。
\end{plainwords}

\begin{definition}[平均奖励]
\[
\rho^\pi = \lim_{T\to\infty} \frac{1}{T} \mathbb{E}_\pi\left[\sum_{t=0}^{T-1} r_t\right]
\]
\end{definition}

\begin{theorem}[平均奖励Bellman方程]
\[
\rho^* + h^*(s) = \max_a \left[R(s,a) + \sum_{s'} P(s'|s,a) h^*(s')\right]
\]
其中$h^*$是差分值函数。
\end{theorem}

\begin{theorem}[R-learning收敛]
R-learning在遍历MDP上收敛到$(\rho^*, h^*)$。
\end{theorem}

%=============================================================================
\section{GPU上的动态规划}
%=============================================================================

\begin{plainwords}
值迭代每个状态更新是\textbf{独立的}——可以并行！
\end{plainwords}

\begin{calculation}[GPU加速]
$|S| = 100000$，值迭代每次遍历所有状态：

\textbf{CPU}：$100000$个状态顺序更新，约40ms

\textbf{GPU}（1000核）：并行更新，约0.4ms

加速100倍！（实际受内存带宽限制，约50倍）
\end{calculation}

\begin{keypoint}[向量化环境]
不仅planning可以并行，\textbf{数据收集}也可以：

同时运行1024个环境，一次收集1024个episode！

这是Isaac Gym的核心思想：GPU原生仿真 + 向量化环境 → 训练速度提升10000倍。
\end{keypoint}

\begin{keypoint}[表格方法的极限]
\begin{center}
\begin{tabular}{l|c|c}
\toprule
问题 & $|S|$ & 可行？ \\
\midrule
井字棋 & $\sim 10^3$ & \checkmark \\
GridWorld 100×100 & $10^4$ & \checkmark \\
Atari（图像） & $\sim 10^{50000}$ & $\times$ \\
围棋 & $\sim 10^{170}$ & $\times$ \\
\bottomrule
\end{tabular}
\end{center}

当$|S| > 10^6$，即使GPU也无法存储值函数表 → 需要\textbf{函数逼近}（下周）
\end{keypoint}

%=============================================================================
\section{收敛性理论大总结}
%=============================================================================

\begin{center}
\begin{tabular}{l|c|c|c}
\toprule
算法 & 收敛性 & 条件 & 速度/界 \\
\midrule
\multicolumn{4}{c}{\textbf{Planning（已知模型）}} \\
\midrule
值迭代 & 全局 & $\gamma < 1$ & $O(\gamma^k)$ \\
策略迭代 & 全局 & 有限MDP & 有限步 \\
\midrule
\multicolumn{4}{c}{\textbf{表格RL}} \\
\midrule
Q-learning & 全局 w.p.1 & 无限访问 & $O(1/k)$ \\
SARSA & 全局 w.p.1 & GLIE & $O(1/k)$ \\
TD($\lambda$) & 全局 & On-policy & 依赖$\lambda$ \\
UCB-VI & Regret界 & Episodic & $\tilde{O}(\sqrt{H^3SAK})$ \\
\midrule
\multicolumn{4}{c}{\textbf{线性函数逼近}} \\
\midrule
Linear TD & 收敛 & On-policy & 放大$\frac{1}{\sqrt{1-\gamma^2}}$ \\
Linear Q & \textbf{可能发散} & Off-policy & — \\
GTD2 & 收敛 & Off-policy & $O(1/k)$ \\
LSTD & 闭式解 & 批量 & $O(\sqrt{d/n})$ \\
\midrule
\multicolumn{4}{c}{\textbf{Multi-Agent}} \\
\midrule
Minimax-Q & 纳什均衡 & 两人零和 & $O(\gamma^k)$ \\
独立Q & \textbf{可能不收敛} & 一般和 & — \\
\bottomrule
\end{tabular}
\end{center}

\begin{magicbox}
\textbf{本讲核心信息}：
\begin{enumerate}
    \item \textbf{MDP = Bandit + 状态转移}，从单步决策到序贯决策
    \item \textbf{Bellman方程}：$V^* = T^* V^*$，收缩映射保证唯一解
    \item \textbf{UCB-VI}：把上周的UCB推广到MDP，Regret $\tilde{O}(\sqrt{H^3SAK})$
    \item \textbf{Linear TD}：投影Bellman方程，误差被$(1-\gamma^2)^{-1/2}$放大
    \item \textbf{Deadly Triad}：函数逼近 + Bootstrapping + Off-policy = 发散风险
    \item \textbf{Simulation Lemma}：模型误差被$(1-\gamma)^{-2}$放大
    \item \textbf{Performance Difference}：用$A^\pi$评估新策略（TRPO/PPO的理论基础）
\end{enumerate}
\end{magicbox}

%=============================================================================
\section*{课后练习}
%=============================================================================

\begin{enumerate}
    \item \textbf{Bellman方程}：3状态MDP（状态1奖励1转到2，状态2奖励0以0.5概率留在2或转到3，状态3奖励2转到1），$\gamma=0.9$。解$V^*$。
    
    \item \textbf{收缩映射}：证明如果$T$是$\gamma$-收缩，则$T^n$是$\gamma^n$-收缩。
    
    \item \textbf{收敛速度}：$\gamma=0.99$时，值迭代需要多少次迭代将误差从100减到0.01？
    
    \item \textbf{Linear TD误差}：$\gamma=0.95$，最佳线性逼近误差为1。TD(0)的误差上界是多少？
    
    \item \textbf{Simulation Lemma}：$\gamma=0.99$，模型误差$\epsilon_P = \epsilon_R = 0.01$，$R_{max}=1$。策略次优性上界是多少？
\end{enumerate}

\subsection*{答案提示}

\textbf{1.} 解方程组：$V(1) = 1 + 0.9V(2)$，$V(2) = 0.9(0.5V(2) + 0.5V(3))$，$V(3) = 2 + 0.9V(1)$。

\textbf{2.} $\|T^n V - T^n V'\| \leq \gamma\|T^{n-1}V - T^{n-1}V'\| \leq \cdots \leq \gamma^n\|V-V'\|$。

\textbf{3.} $0.99^k < 0.0001$，$k > \frac{\ln 10^6}{\ln(1/0.99)} = \frac{13.8}{0.01} = 1380$次。

\textbf{4.} $\frac{1}{\sqrt{1-0.95^2}} = \frac{1}{\sqrt{0.0975}} = 3.2$。误差上界$\leq 3.2$。

\textbf{5.} $\frac{2}{0.01^2}(0.01 + 0.99 \times 0.01 \times 100) = 20000 \times 1.01 \approx 20200$。

\end{document}